% Blueprint content for leslie2.
% Shares its macros and its figures (figures/) with the full edition
% (content-full.tex), and carries its own connective prose.
% Each edition has its own node bodies, nodes-min/ here and nodes/ there, and
% scripts/check-node-sync.py keeps the frontmatter of each pair in sync.
% Each section says what its objects are for and points at the Lean source; the
% statements are the nodes and the full implementation is the Lean.

\chapter{The framework}

\section{Systems, executions, and trace distributions}

A PLTS is the single model every system of this development is written in.
Executions record whole runs and traces record what is observable of them, and that separation is what lets an implementation be compared with its specification.
A refinement claim here is always an inclusion $\operatorname{ATD}(sys_C) \subseteq \operatorname{ATD}(sys_A)$ of achievable trace distributions, and safety properties are read off the supports of those distributions.
Throughout, $\tau$ is the distinguished silent label of the label type, a label is internal when it equals $\tau$ and external otherwise, $\delta_s$ is the Dirac distribution on $s$, $\mu \triangleright g$ is the pushforward of $\mu$ along $g$, and $s[x := v]$ is $s$ with the field $x$ set to $v$.
See \leanmodule{Leslie2/Systems/System} and \leanmodule{Leslie2/Systems/Trace}.

\input{nodes-min/def-system}
\input{nodes-min/def-exec}
\input{nodes-min/def-trace}
\begin{definition}[Probabilistic execution and its path measure]
  \label{def:probabilistic-execution}
  \lean{PLTS.ProbabilisticExecution}\leanok
  \uses{def:system, def:exec}
  A probabilistic execution is an initial distribution together with a scheduler, and it gives each finite execution the mass of the path it follows: \leandecl{PLTS.ProbabilisticExecution}.
\end{definition}

\input{nodes-min/def-istight}
\begin{definition}[Trace probability]
  \label{def:traceprob}
  \lean{PLTS.System.traceProb}\leanok
  \uses{def:probabilistic-execution, def:trace, def:istight}
  The trace probability of a finite trace under a probabilistic execution is the summed path mass of the tight finite executions carrying that trace: \leandecl{PLTS.System.traceProb}.
\end{definition}

\input{nodes-min/def-achievable}

\section{Weak transitions and the two constructions}

Weak transitions absorb internal computation, and all three of them move a distribution over states rather than a state and run until an internal scheduler halts.
The two constructions closing the section turn those notions into ordinary systems: the weak closure $\cdot^w$ makes a weak transition a single step, and the distribution-monad lift $\mathcal{D}(\cdot)$ makes a distribution over states a single state.
The weak transitions are in \leanmodule{Leslie2/Weak/WeakTransition} and \leanmodule{Leslie2/DistMonad/HyperStep}, and the two constructions in \leanmodule{Leslie2/WeakClosure/WeakClosure} and \leanmodule{Leslie2/DistMonad/DistMonad}.

\input{nodes-min/def-weaktau}
\input{nodes-min/def-hyperstep}
\input{nodes-min/def-weakstep}
\input{nodes-min/def-weakclosure}
\input{nodes-min/def-dist}

\section{Probabilistic simulations}

A simulation of a concrete system by an abstract one is carried by a relation together with a step-matching condition, phrased through a coupling of the two successor distributions.
The two notions of this section differ only in what the abstract system may match a concrete step with: a single step for strong simulation, and a weak transition out of a belief distribution over abstract states for forward simulation.
Every refinement of the case study is proved in the second of them, whose belief distribution is where an abstract commitment is deferred to.
See \leanmodule{Leslie2/Simulation/SimDefs}.

\begin{definition}[Coupling of distributions]
  \label{def:coupling}
  \lean{PLTS.PMFRel}\leanok
  A coupling of two distributions along a relation is a distribution on the product with those two marginals whose support lies inside the relation: \leandecl{PLTS.PMFRel}.
\end{definition}

\begin{definition}[Strong probabilistic simulation]
  \label{def:strongsim}
  \lean{PLTS.StrongProbabilisticSimulation}\leanok
  \uses{def:system, def:coupling}
  A strong probabilistic simulation is a relation on states, holding at the two initial states, in which every concrete step is matched by a single abstract step on the same label whose successor distributions are coupled: \leandecl{PLTS.StrongProbabilisticSimulation}.
\end{definition}

\begin{definition}[Probabilistic forward simulation]
  \label{def:forwardsim}
  \lean{PLTS.ProbabilisticForwardSimulation}\leanok
  \uses{def:weaktau, def:weakstep, def:coupling}
  A probabilistic forward simulation, the notion of Segala~\cite{Segala95}, relates each concrete state to a distribution over abstract states and matches every concrete step by a weak transition out of that distribution under a coupling: \leandecl{PLTS.ProbabilisticForwardSimulation}.
\end{definition}


\section{Building blocks: functional simulation and the constructions}

Functional simulation is the one lemma the soundness proof is built on: a label-preserving map on states whose pushforward matches every step, and which is therefore checkable with no coupling exhibited.
See \leanmodule{Leslie2/Simulation/TraceMap}, \leanmodule{Leslie2/Simulation/StrongSoundness}, \leanmodule{Leslie2/WeakClosure/TraceProb} and \leanmodule{Leslie2/DistMonad/DistTrace}.

\input{nodes-min/lem-map}
\input{nodes-min/thm-strongsound}
\input{nodes-min/thm-weakclosure-eq}
\input{nodes-min/thm-dist-eq}

\section{Soundness of probabilistic forward simulation}
\label{sec:soundness}

Forward simulation is definitionally a strong simulation into $\mathcal{D}(sys_A^w)$, and soundness is that observation together with the two preservation results above.
See \leanmodule{Leslie2/Results}.

\subsection{The transition-level correspondence}

A step of $\mathcal{D}(sys_A^w)$ out of a distribution $\mu$ is a hyper-step of $sys_A^w$, and it has to be identified with a single weak transition of $sys_A$ out of $\mu$.
Source decomposition, source mixing and target convexity are the three families that identification rests on, and the correspondence lemma is their consequence.
See \leanmodule{Leslie2/Weak/WeakChar} and \leanmodule{Leslie2/Simulation/Equivalences}.

\input{nodes-min/lem-decomp}
\input{nodes-min/lem-mix}
\input{nodes-min/lem-bind}
\begin{lemma}[Hyper-step as weak transition]
  \label{lem:hyperstep-weakstep}
  \lean{PLTS.hyperStep_weakClosure_of_weakTau, PLTS.weakTau_of_hyperStep_weakClosure, PLTS.hyperStep_weakClosure_of_weakStep, PLTS.weakStep_of_hyperStep_weakClosure}\leanok
  \uses{def:hyperstep, def:weakclosure, lem:decomp, lem:mix, lem:bind}
  A hyper-step of $sys^w$ out of a distribution is exactly a weak transition of $sys$ out of that distribution, silent or external according to the label: \leandecl{PLTS.hyperStep_weakClosure_of_weakTau}, \leandecl{PLTS.weakTau_of_hyperStep_weakClosure}, \leandecl{PLTS.hyperStep_weakClosure_of_weakStep}, \leandecl{PLTS.weakStep_of_hyperStep_weakClosure}.
\end{lemma}

\begin{proof}\leanok
  \uses{lem:decomp, lem:mix, lem:bind}
  See the Lean proofs of \leandecl{PLTS.weakTau_of_hyperStep_weakClosure} and \leandecl{PLTS.weakStep_of_hyperStep_weakClosure} and of their two converses.
\end{proof}


\subsection{The reduction}

The obligation a refinement proof is left with is one relation checked one concrete transition at a time.

\begin{theorem}[Forward simulation $=$ strong simulation into $\mathcal{D}(\cdot^w)$]
  \label{thm:forward-equiv}
  \lean{PLTS.probabilisticForwardSimulation_iff_strong_dist_weakClosure}\leanok
  \uses{def:forwardsim, def:strongsim, def:dist, def:weakclosure, lem:hyperstep-weakstep}
  A relation is a forward simulation of $sys_C$ by $sys_A$ exactly when it is a strong simulation of $sys_C$ by $\mathcal{D}(sys_A^w)$: \leandecl{PLTS.probabilisticForwardSimulation_iff_strong_dist_weakClosure}.
\end{theorem}

\begin{proof}\leanok
  \uses{lem:hyperstep-weakstep}
  See the Lean proof of \leandecl{PLTS.probabilisticForwardSimulation_iff_strong_dist_weakClosure}.
\end{proof}

\input{nodes-min/thm-forward-sound}

\subsection{Consequences and constructions}

Safety properties are stated on the supports of achievable trace distributions, and the trace-support transfer below carries such a property from a system to every system it simulates.
The transfer is in \leanmodule{Leslie2Protocols/Framework/TraceDistributionSupport} and the composition operators are in \leanmodule{Leslie2/ProcessAlgebra/Composition} and \leanmodule{Leslie2/ProcessAlgebra/Abstract}.
The synchronised product of the components is in \leanmodule{Leslie2Protocols/Framework/SynchronisedProduct} and the idle-padded instance families are in \leanmodule{Leslie2Protocols/Framework/LoopsAndInstanceFamilies}.
The product of a family of components each read along its own pullback is in \leanmodule{Leslie2Protocols/Framework/SynchronisedProductAlongPullbacks}.
The two couplings a Dirac lift admits --- a Dirac source matched by a related target, and a source pushed forward along a map matched outcome by outcome --- are in \leanmodule{Leslie2Protocols/Framework/DiracRelationCoupling}.
The restriction along the left embedding and the precongruence it enjoys are both in \leanmodule{Leslie2Protocols/Framework/Relabel}.
The read-back congruence, in \leanmodule{Leslie2Protocols/Framework/FinerAlphabetCongruence}, carries a simulation over to a finer alphabet, and \leanmodule{Leslie2Protocols/Framework/WeakTransitionsFromChains} supplies the weak runs of unbounded length --- a chain of silent steps, closed or not by one external step --- that the simulations of the case study give as matching runs.
The last statement of the list is the same congruence property at the level of labelled transition systems, for every operator at once, which is what a substitution inside a composition of components rests on; it is in \leanmodule{Leslie2Protocols/Framework/Congruence}.

\begin{theorem}[Trace-support transfer]
  \label{thm:trace-support}
  \lean{PLTS.exists_exec_of_traceProb_ne_zero}\lean{PLTS.is_exec_of_probOf_ne_zero}\lean{PLTS.safety_transfer}\leanok
  A trace of positive probability under a Dirac-initialised probabilistic execution comes from a genuine execution whose events characterise membership in that trace, so safety of trace supports transfers along inclusion of achievable trace distributions: \leandecl{PLTS.exists_exec_of_traceProb_ne_zero}, \leandecl{PLTS.is_exec_of_probOf_ne_zero}, \leandecl{PLTS.safety_transfer}.
\end{theorem}

\begin{proof}\leanok
  \uses{def:probexec, def:traceprob, def:exec, def:achievable}
  See the Lean proofs of \leandecl{PLTS.exists_exec_of_traceProb_ne_zero} and \leandecl{PLTS.safety_transfer}.
\end{proof}

\input{nodes-min/def-process-algebra}
\begin{definition}[Extended alphabets and restriction]
  \label{def:relabel}
  \lean{PLTS.instSilentSum}\lean{PLTS.System.relabel}\lean{PLTS.System.relabel_isLTS}\lean{PLTS.System.abstract_isLTS}\leanok
  A composition whose components synchronise on labels it must not expose is built over the extended alphabet $\mathrm{Label} \oplus \mathrm{Extra}$, whose auxiliary labels are observable and hence hideable, and restriction along the left embedding retains the $\mathrm{inl}$-labelled transitions alone: \leandecl{PLTS.instSilentSum} and \leandecl{PLTS.System.relabel}.
\end{definition}

\begin{definition}[Idle padding and instance families]
  \label{def:idle-family}
  \lean{PLTS.System.withIdle}\lean{PLTS.System.family}\leanok
  \uses{def:process-algebra}
  Idle self-loops on the labels a component does not own make full synchronisation emulate sync-set composition, and the $\mathbb{N}$-indexed instance family moves one instance per owned or silent label, broadcasts corruption to all of them, and idles otherwise: \leandecl{PLTS.System.withIdle} and \leandecl{PLTS.System.family}.
\end{definition}

\begin{definition}[The full-synchronisation product]
  \label{def:synchronised-product}
  \lean{PLTS.System.synchronisedProduct}\lean{PLTS.System.synchronisedProduct_isLTS}\leanok
  \uses{def:idle-family}
  The full-synchronisation product steps every component of a finite family at once on a visible label and exactly one component on $\tau$, ownership of a label being carried by the components and not by the operator: \leandecl{PLTS.System.synchronisedProduct}.
\end{definition}

\input{nodes-min/def-strongfunctional}
\begin{theorem}[Transitivity of probabilistic forward simulation]
  \label{thm:forward-trans}
  \lean{PLTS.ProbabilisticForwardSimulation.trans}\leanok
  \uses{def:forwardsim}
  Forward simulations compose, along the composite of the two relations: \leandecl{PLTS.ProbabilisticForwardSimulation.trans}.
\end{theorem}

\begin{proof}\leanok
  \uses{def:forwardsim, def:weakstep, lem:hyperstep-weakstep}
  See the Lean proof of \leandecl{PLTS.ProbabilisticForwardSimulation.trans}.
\end{proof}

\input{nodes-min/thm-precong-parallel}
\input{nodes-min/thm-precong-interleave}
\input{nodes-min/thm-precong-abstract}
\input{nodes-min/thm-precong-relabel}
\begin{theorem}[Read-back congruence]
  \label{thm:mapidle-cong}
  \lean{PLTS.ForwardSimulation.mapIdle}\leanok
  \uses{def:idle-family, def:forwardsim}
  Forward simulation is a congruence for reading both systems over a finer alphabet along one partial label map: delegated labels are matched through the simulation and unmapped labels by idling, under the sole hypothesis that the silent labels of the two alphabets correspond and nothing else delegates to $\tau$: \leandecl{PLTS.ForwardSimulation.mapIdle}.
\end{theorem}

\begin{proof}\leanok
  \uses{def:idle-family, def:forwardsim}
  See the Lean proof of \leandecl{PLTS.ForwardSimulation.mapIdle}.
\end{proof}

\begin{theorem}[Forward simulation is a congruence for the composition operators]
  \label{thm:forwardsim-congruence}
  \lean{PLTS.ForwardSimulation.parallel_right}\lean{PLTS.ForwardSimulation.parallel_left}\lean{PLTS.ForwardSimulation.synchronisedProduct}\lean{PLTS.ForwardSimulation.abstract}\lean{PLTS.ForwardSimulation.relabel}\lean{PLTS.ForwardSimulation.trans}\leanok
  \uses{def:forwardsim, def:process-algebra, def:relabel, def:syncproduct}
  Forward simulation between labelled transition systems survives each operator a composition is built from --- parallel composition on either side, the full-synchronisation product of a finite family, hiding a label set, and restriction along the left embedding --- and two forward simulations compose: \leandecl{PLTS.ForwardSimulation.parallel_right}, \leandecl{PLTS.ForwardSimulation.parallel_left}, \leandecl{PLTS.ForwardSimulation.synchronisedProduct}, \leandecl{PLTS.ForwardSimulation.abstract}, \leandecl{PLTS.ForwardSimulation.relabel} and \leandecl{PLTS.ForwardSimulation.trans}.
  Replacing one component of a composition by a system that simulates it is an application of these, and every substitution inside a gather instance and inside a graded-agreement round is proved that way.
\end{theorem}

\begin{proof}\leanok
  \uses{def:forwardsim, def:process-algebra, def:syncproduct, def:weakstep}
  See the Lean proofs of \leandecl{PLTS.ForwardSimulation.parallel_right}, \leandecl{PLTS.ForwardSimulation.synchronisedProduct}, \leandecl{PLTS.ForwardSimulation.abstract}, \leandecl{PLTS.ForwardSimulation.relabel} and \leandecl{PLTS.ForwardSimulation.trans}.
\end{proof}


A system is often written with a state component that no transition's firing depends on, carried so that an invariant can be stated over it, and read at most to decide a value a step carries on its label, the \emph{announced} value, and removing such a component must leave the observable behaviour untouched.
A component of that kind is an auxiliary variable in the sense of Abadi and Lamport~\cite{AL91}.
The definition below states its removal for systems whose steps are distributions, and the equality of achievable trace distributions it yields rests on Theorem~\ref{thm:forward-sound} alone.
Congruences carry a removal through parallel composition, abstraction and restriction, so it composes as a simulation does.
The construction is in \leanmodule{Leslie2Protocols/Framework/AuxiliaryVariableRemoval}.

\begin{definition}[Auxiliary-variable removal]
  \label{def:auxiliary-variable-removal}
  \lean{PLTS.AuxiliaryVariableRemoval}\lean{PLTS.System.LabelSaturated}\lean{PLTS.SeparatesSilent}\lean{PLTS.AuxiliaryVariableRemoval.achievableTraceDists_eq}\leanok
  \uses{def:forwardsim, def:achievable, def:strongfunctional, thm:forward-sound}
  An auxiliary-variable removal of $sys_A$ onto $sys_0$ is a map $\pi$ on states together with an identification $\varphi$ of labels, under which the two systems match each other's transitions along $\pi$ and $\varphi$ and $\tau$ is the one label with the $\varphi$-image of $\tau$: \leandecl{PLTS.AuxiliaryVariableRemoval}, \leandecl{PLTS.SeparatesSilent}.
  At $\varphi = \mathrm{id}$ the two systems have the same achievable trace distributions, a general $\varphi$ reaches that equality once every label it identifies with a different label is hidden, and a removal survives parallel composition, abstraction and restriction against a component that is label-saturated for $\varphi$: \leandecl{PLTS.AuxiliaryVariableRemoval.achievableTraceDists_eq}, \leandecl{PLTS.System.LabelSaturated}.
\end{definition}


\chapter{The ABA case study}

The setting is $n$ processes, each holding an input bit, of which at most $f$ are Byzantine, with $3f < n$, communicating by asynchronous message passing.
The development is in the opt-in Lake library \leanmodule{Leslie2Protocols}.

\section{Setting and method}

Three protocols occur in the case study, each given as properties over the calls and returns of its interface that quantify over processes never corrupted.
What this development establishes is Validity and Agreement for ABA, machine-checked; Termination, Unpredictability and fairness are out of scope, and the weak common coin is held at specification level, so $\varepsilon$-Goodness enters as an assumption.

Asynchronous Byzantine Agreement (ABA) takes a bit from each process and may return a bit.
\begin{itemize}
  \item Validity: a correct process returns only a bit a correct process had as input.
  \item Agreement: two correct returns carry the same bit.
  \item $\varepsilon$-sure Termination: once $n-f$ correct processes take part, all correct processes return with probability $1-\varepsilon$.
\end{itemize}

Graded Binding Crusader Agreement (GBCA) takes a bit and may return a graded value in $\{(0,2),\ (0,1),\ (\bot,0),\ (1,1),\ (1,2)\}$.
\begin{itemize}
  \item Validity: a correct process failing to return $(b,2)$ implies a correct process had $1-b$ as input.
  \item Graded Agreement: returns $(b,X)$ and $(b',X')$ of correct processes satisfy $b = 0 \Rightarrow b' \neq 1$ and $X = 2 \Rightarrow X' \neq \bot$.
  \item Binding: when a correct process returns, a bound value $b$ is defined and no extension of the execution returns $(1-b,1)$ or $(1-b,2)$ to a correct process.
  \item Termination: once $n-f$ correct processes take part, all correct processes return.
\end{itemize}

Weak Common Coin (WCC) takes no input and may return a bit.
\begin{itemize}
  \item $\varepsilon$-Goodness: with probability $\varepsilon$ all correct processes return $0$, and with probability $\varepsilon$ all return $1$.
  \item Unpredictability: all outcomes remain possible until $f+1$ processes have called.
  \item $(1-\failmass)$-sure Termination: once $n-f$ correct processes take part, all correct processes return with probability $1-\failmass$.
\end{itemize}

Every protocol with sub-protocols is treated in two moves: substitution, in which each sub-implementation is exchanged for the corresponding specification in a fixed context and which is sound by precongruence, and the core simulation, in which the core algorithm beside the specifications of its sub-protocols is simulated by the protocol's own specification.
Composing the two inclusions places the implementation inside the specification, and the recursion stops at WCC, which is left at specification level.
GBCA is verified against two implementations: a direct one, which calls no sub-protocol, and a gather-based one, through which the recursion continues into gather and reliable broadcast.

\paragraph{Object names.}
Transition System~1 is the ABA specification, Transition System~2 the GBCA specification, Transition System~3 the WCC specification, Transition System~4 the gather specification and Transition System~6 the BRB specification; the numbering is the source's, whose Transition Systems~5 and~7 (SRSD and AVSS) are not encoded, and the names are used throughout this section and in every node of it.
A transition is cited by its Lean constructor name, a deviation label in parentheses marks a point where an encoding departs from~\cite{LeslieBP}, and where a node names ABDY22's Algorithm~$k$ it means the numbering of~\cite{ABDY22}.
A \emph{ghost output} is a value a specification holds as state and announces on a return label, which no program reads and on which no transition's firing depends (D29); the round's bound bit and a gather instance's core are the two.
An \emph{auxiliary-variable removal} deletes a state component that no transition's firing depends on, and the \emph{ghost-free system} of a protocol is its implementation with the adversary's ghost deleted.
The grades are $2$, $1$ and $0$ here and in~\cite{ABDY22} and~\cite{AFW25}, where Transition System~2 of the source blueprint writes $A$, $B$ and $C$.
Bracha's third message level is $\mathrm{VOTE}$ here and in the source blueprint, where both~\cite{ABDY22} and~\cite{AFW25} write $\mathrm{READY}$.
Its steps fire on the received messages and at the quorums of Algorithm~1 of~\cite{AFW25}, where Algorithm~6 of the source blueprint fires the echo step on a received $\mathrm{INIT}$ alone and puts $n-f$ at every quorum.
The naming difference carries no deviation label, and the rest is D34.

\paragraph{Outside the model.}
Two features of~\cite{LeslieBP} are deliberately not encoded: fairness, whose markings bear only on liveness, and partial information, whose observation functions and belief construction are not needed because every adversary here is omniscient.

\paragraph{Notation.}
Write $[n]$ for the identifier set and $F \subseteq [n]$ for the corrupted set.

\section{Deviations from the source blueprint}
\label{sec:aba-deviations}

% The registry of deviations from the source blueprint, default edition.
% The full edition's registry is deviations-full.tex.
% The overview table is written by hand, and scripts/check-deviations-sync.py
% checks it against the entries below and against the full edition.
Each departure from~\cite{LeslieBP} carries a label, cited at the point where it applies.
The table lists the labels in use with a headline each, and the entries below state, under each label, the departure and its reason.
\begin{center}
\footnotesize \setlength{\tabcolsep}{4pt}
\begin{tabular}{lp{12cm}}
\hline
Label & Headline \\
\hline
D1 & corruption as a guarded transform, broadcast on fail \\
D4 & no WCC \texttt{guess} label or state field \\
D5 & a network component holding per-sender sent sets \\
D8 & round sends and returns guarded by the process's own input \\
D9 & $0$-based round indexing \\
D11 & Byzantine call and return transitions at each sub-protocol interface \\
D12$'$ & per-process DECIDED sent and received sets, with equivocation \\
D13 & a ghost $input$ and an $f+1$ support guard on the decision \\
D14 & $f+1$ support guards on the GBCA specification's binding and returns \\
D15 & $F$-blind support counts \\
D16 & a two-phase abstract state deciding at the first visible return \\
D17 & a fourth coin outcome $\dagger$ of mass $\failmass$ \\
D18 & ABDY22's Algorithm~6 transcribed in full, ECHO5 level included \\
D19 & negative binding, the exclusion set as the state \\
D21 & Transition System~1 as eight transitions over a mode \\
D22 & retention with termination \\
D23 & a corrupted process's program replaced by a self-loop \\
D24 & the gather-based round's grade read off the second gather's counts \\
D26 & gather calls and committed entries in separate fields \\
D27 & the BRB leader's recorded input split from the committed value \\
D29 & ghost outputs carried down as auxiliary variables \\
D30 & a parametric per-round ghost held by the network \\
D31 & the coin's threshold counting callers alone \\
D32 & each BIND payload sent by reliable broadcast \\
D33 & both gathers of the round in the binding form \\
D34 & Bracha's rounds as published, three echo triggers and the $(n+f)/2$ quorum \\
D35 & the gather transitions carrying the order of the sends \\
D36 & the input-enabledness loop guarded on its variables \\
\hline
\end{tabular}
\end{center}
\begin{deviations}
\deviation{D1}{corruption as a guarded transform, broadcast on fail}
\departure{Corruption is a guarded deterministic transform $\mathrm{corrupt}$, which inserts $id$ into $F$ and writes nothing else, broadcast to every component on $\mathrm{fail}(id)$ under the guard $id \notin F \wedge |F| < f$ at the component that owns the corrupted set, its half at the process the label names being the replacement of D23.}
\reason{The budget is a guard of the one component that owns the corrupted set, so no headline carries a hypothesis on the trace, and the broadcast keeps every component's copy of the set equal, no component reading another's state.}
\deviation{D4}{no WCC \texttt{guess} label or state field}
\departure{The WCC \texttt{guess} label and the \texttt{guess} state field are omitted.}
\reason{The label and the field exist solely for Unpredictability, which is out of scope.}
\deviation{D5}{a network component holding per-sender sent sets}
\departure{The network holds the set $sent(r,j)$ of the round-$r$ messages process $j$ has multicast; a multicast inserts into it under guards at the sender, a delivery checks membership and consumes nothing, a threshold counts distinct senders, and arbitrary content enters only through the injection under $k \in F$.}
\reason{The source carries no network, and Section~2 of~\cite{ABDY22} and Section~2.1 of~\cite{AFW25} state it in one sentence each, links that do not drop messages and reliable point-to-point messages with arbitrary content from a faulty process.
  The sets absorb duplication, re-delivery and per-link order, and they admit the equivocation that D12$'$ mirrors for DECIDED.}
\deviation{D8}{round sends and returns guarded by the process's own input}
\departure{Every protocol send of a round, the round's three graded returns, the ABA return and the DECIDED relay are guarded by the acting process's own input, which is the reading of~\cite{ABDY22}.}
\reason{There every party receives an input (Definition~2.2, Definition~3.2 stating Validity over the inputs of the non-faulty parties), and participation begins with the input send: ``To participate in the algorithm, a party $p$ begins by sending its input value $x$ to all parties'' (Section~4.1), ``a party starts by sending its input value $x$ to all other parties'' (Section~4.2), and Algorithm~6 is ``nearly identical to Algorithm~4 up through line~17''.
  The uncalled process and the caller quorum are Transition System~2's, and the guard is what discharges that quorum: the correct senders of an $n-f$ quorum of received $\mathrm{INPUT}$ messages are callers.
  Without the guard the source's Algorithm~2 does not refine the source's Transition System~2; the guard holds across the gather layer of both chains, and Bracha's rounds carry none, since only the leader is called.}
\deviation{D9}{$0$-based round indexing}
\departure{$0$-based round indexing.}
\reason{The round counter and every round-indexed family range over $\mathbb{N}$, whose first element indexes the initial round variables.}
\deviation{D11}{Byzantine call and return transitions at each sub-protocol interface}
\departure{Byzantine call and return transitions at each sub-protocol interface.
  In an implementation the graded-agreement call and return have no transition at the process their label names: \leandeclnamed{byzantineCallGIdle}{PLTS.ABA.Implementation.ProgramStep.byzantineCallGIdle} and \leandeclnamed{byzantineRetGIdle}{PLTS.ABA.Implementation.ProgramStep.byzantineRetGIdle} carry $k \neq j$, and \leandeclnamed{corruptedIdle}{PLTS.ABA.Implementation.ProgramStep.corruptedIdle} excludes the labels of $\mathrm{actsAt}\,j$.}
\reason{An implementation therefore neither calls into nor returns from a round at a corrupted process, both labels being hidden, and the composed system answers them, so an implementation is included in its composed system and the inclusion holds in that direction alone.
  The authorisation $k \in F$ belongs to the network that surrounds the instance, where the call or return label stays visible, so the instance carries the effect of a Byzantine transition and never its guard, and a replaced program (D23) takes no transition of its own on those labels.}
\deviation{D12$'$}{per-process DECIDED sent and received sets, with equivocation}
\departure{Per-process DECIDED sent sets and received sets, delivery per (receiver, sender, bit), and Byzantine equivocation in those sent sets.}
\reason{A single DECIDED entry per process cannot send $\mathrm{DECIDED}\,0$ to one process and $\mathrm{DECIDED}\,1$ to another, an under-approximation inconsistent with the equivocating sent sets of D5, and delivery per (receiver, sender, bit) is what splits each delivery into the network's half and the receiver's.}
\deviation{D13}{a ghost $input$ and an $f+1$ support guard on the decision}
\departure{What Validity is proved from on Transition System~1, namely a ghost $input$ of the genuine $\mathrm{callABA}$ events and an $f+1$ support guard \leandecl{PLTS.ABA.InputSupport} on the decision.
  The ghost has one correct writer, the interface transition \leandeclnamed{callSet}{PLTS.ABA.SpecStep.callSet}, whose guard is the empty entry, so its write is a first write.
  The input-enabledness loop \leandeclnamed{callLoop}{PLTS.ABA.SpecStep.callLoop} carries the call label at a filled entry and writes nothing (D36).}
\reason{Without the guard the specification violates Validity: at $n = 4$, $f = 1$, with the recorded inputs $1,0,0,0$ and the sole holder of $1$ corrupted, a decision of $1$ has every never-corrupted process inputting $0$.
  Among $f+1$ supporters some process is never corrupted, since at most $f$ ever are, so the decision is attributed to a never-corrupted caller with no guard that reads the future.}
\deviation{D14}{$f+1$ support guards on the GBCA specification's binding and returns}
\departure{$f+1$ support guards on the GBCA specification's binding, on the dissenting bit at its grade-1 return, and on each of the two bits at its grade-0 return.}
\reason{Transition System~2's singular witness may be corrupted later in the trace, after which nothing attributes the outcome to a never-corrupted input, and $\mathrm{hybrid}$ over that form violates Validity: at $n = 4$, $f = 1$ with inputs $1,0,0,0$, binding at $1$ off the sole holder of $1$, grade-1 adoption everywhere, unanimity in the next round and the corruption of that holder return $1$ where every never-corrupted process input $0$.}
\deviation{D15}{$F$-blind support counts}
\departure{The $F$-blind form of those support counts, namely $f+1$ distinct callers-or-corrupted.}
\reason{The count is monotone in $F$ and in the calls, so no later corruption invalidates it, and among $f+1$ supporters some member lies outside the final $F$, which recovers a never-corrupted genuine caller; filling the entries of corrupted processes at the specification instead would require knowing which process is corrupted later, a prophecy no forward simulation has.}
\deviation{D16}{a two-phase abstract state deciding at the first visible return}
\departure{The two-phase abstract state used by the core simulation, which defers the specification's $\tau$-progress to a single \leandeclnamed{decide}{PLTS.ABA.SpecStep.decide} step under the first visible return.}
\reason{An abstract state that decides when a round excludes a bit can be contradicted afterwards: a late caller of the other bit enables a grade-0 return at that round, the next round spares the opposite bit and is graded~2, and the decision falls against the committed value.
  The abstract state therefore commits as late as a forward simulation permits.}
\deviation{D17}{a fourth coin outcome $\dagger$ of mass $\failmass$}
\departure{A fourth outcome $\dagger$, of mass $\failmass$, on the single distribution both coin resolutions follow, where the source's two range over $\{0,1,\top\}$.
  At Transition System~3 it delivers to no process, and at Transition System~1 it puts the system into the mode $noTransitionEnabled$, which enables no silent transition.}
\reason{The implementation the coin specification stands for reconstructs the coin from an $\varepsilon$-correct verifiable secret sharing whose delivery holds only up to a failure probability, and the specification carries that failure as an outcome of its own rather than as a side condition.}
\deviation{D18}{ABDY22's Algorithm~6 transcribed in full, ECHO5 level included}
\departure{The GBCA implementation transcribes ABDY22's Algorithm~6 in full, over the message levels INPUT, ECHO, VOTE, BIND and ECHO5, in place of the source's compression of it.}
\reason{That compression elides the ECHO5 level and takes $f{+}1$ received $\mathrm{VOTE}\,v$ messages as the grade-$1$ witness where Algorithm~6 takes $f{+}1$ received $\mathrm{BIND}\,v$ messages, which violates Binding.}
\deviation{D19}{negative binding, the exclusion set as the state}
\departure{Binding is negative and the exclusion set is the state, the GBCA specification carrying $excluded \subseteq \{0,1\}$ under one inserting internal transition and no removing transition, in place of the source's bound value $bind \in \{0,1,\bot\}$, so that a return's guard on the bit it announces is $1-bnd \in excluded$ where the source's is $bind = bnd$ (Definition~\ref{def:aba-gbca-spec}).}
\reason{With the exclusion set, Graded Agreement is the guard pair $v \notin excluded \wedge 1-v \in excluded$ and Binding is the guard on the announced bit, both consequences of the set's monotonicity alone, with no auxiliary invariant and no quorum arithmetic, and the once-only guard on the exclusion is what keeps a round completable after a grade-2 return.}
\deviation{D21}{Transition System~1 as eight transitions over a mode}
\departure{Transition System~1 is eight transitions over the state $\{input, ret, F, val, mode\}$, in place of the source's per-process $call$ field and its four-valued $coin \in \{0,1,\bot,\top\}$.}
\reason{Transition System~1's re-proposal transition fires with $val$ already written, so its unanimity transition can overwrite the decision between two returns, which violates Agreement; with \leandeclnamed{decide}{PLTS.ABA.SpecStep.decide} the sole writer of $val$ under $val = \bot$, the decision is written once by the shape of the system.
  The flip names no bit because a bit entering through the one probabilistic transition would re-admit a Validity-breaking decision at probability $\varepsilon$; the identification of the flip's outcome with the coin agreeing with a round's value belongs to a refinement's outcome coupling.}
\deviation{D22}{retention with termination}
\departure{Retention with termination, a process holding its variables in every round it has touched in the finite map of \leandecl{PLTS.ABA.Implementation.RoundVariablesMap}, where a round never touched reads as the initial variables, and the round's message transitions reading the variables their own label tags rather than the one the round loop stands at.
  Termination is \leandeclnamed{terminate}{PLTS.ABA.Implementation.ProgramStep.terminate}, which fires once the process's own return has fired and $2f+1$ received DECIDED messages stand in its variables, sets the flag alone, and carries the replacement guard of D23.}
\reason{ABDY22 separates deciding from terminating, and its amplification argument (Lemmas~4.6 and~E.5) counts the echoes a decided process keeps sending, so a process must answer the messages of every round it has touched; the finite map keeps that state finite by construction, and $2f+1$ received DECIDED messages stand behind $f+1$ correct senders, the relay threshold, so a terminated process holds no other back.}
\deviation{D23}{a corrupted process's program replaced by a self-loop}
\departure{When a process is corrupted its program is replaced, \leandeclnamed{failSelf}{PLTS.ABA.Implementation.ProgramStep.failSelf} writing the flag \leandecl{PLTS.ABA.RoundLoopVariables.corrupted}, every correct transition requiring an unreplaced program, and the replaced program being the single self-loop \leandeclnamed{corruptedIdle}{PLTS.ABA.Implementation.ProgramStep.corruptedIdle}.
  Both networks carry the Byzantine visible return $\mathrm{retABA}(id,b)$ under the authorisation $id \in F$ alone, and Transition System~1's image of the replacement is \leandeclnamed{callByzantine}{PLTS.ABA.SpecStep.callByzantine}, recording a bit unrelated to the one its label declares, beside \leandeclnamed{retByzantine}{PLTS.ABA.SpecStep.retByzantine}, returning any bit at any time without moving the state.}
\reason{A corrupted process is the adversary's, so the model lets it call one bit and record another and return either bit at any time; the replacement flag is the one thing a program reads about corruption, so no program reads the corrupted set or the budget, and the two corrupted transitions of Transition System~1 add behaviour and remove none, which is why both trace predicates are read at never-corrupted returners.}
\deviation{D24}{the gather-based round's grade read off the second gather's counts}
\departure{The gather-based GBCA implementation follows Algorithm~4 of~\cite{AFW25} at $R = 2$, with the grade read off the second gather's counts in place of the approximate-agreement subroutine of lines~7 and~8: $(v, 2)$ at $|{\rm dom}\, h| - f$ entries of $v$, $(v, 1)$ at $f{+}1$, $(\bot, 0)$ below both.}
\reason{The source supplies no gather-based GBCA, its gather serving the coin construction, which stays at specification level here.
  The local count carries the exclusivity the subroutine establishes, so its message exchange and its thresholds add a step the round's safety does not need.}
\deviation{D26}{gather calls and committed entries in separate fields}
\departure{The gather specification records calls and committed entries in separate fields, a per-entry internal commit writing $val(k)$ once under $k \in F \vee call(k) = v$, in place of Transition System~4's transition writing $call$ itself at a corrupted process.
  The transition writing the core fires on a set of committed entries of size at least $n-f$, where Transition System~4 fires on more than $n-f$ recorded calls and puts no size bound on the set it takes.}
\reason{Entries reach the gather by reliable broadcast, so a process corrupted after a correct call may still direct its committed entry until first use, and a specification that fixed the entry at the call would refuse that execution; the size guard is what Binding Common Core states and Transition System~4's transition, which admits the empty set, does not enforce.}
\deviation{D27}{the BRB leader's recorded input split from the committed value}
\departure{The BRB specification splits the leader's recorded input from the committed value, an internal commit writing the value once under $\mathit{ldr} \in F \vee input = m$, in place of Transition System~6's single call field and its corrupted-leader transition, which only bars the instance from returning.}
\reason{Under Transition System~6 a correct call fixes every return, where Bracha's rounds with the leader corrupted after its $\mathrm{INIT}$ can deliver another value until some correct process holds an $\mathrm{ECHO}$ quorum, so that specification would exclude its own implementation under D1; the split closes the window at the commit, a leader corrupted after its call may commit any value, and Validity holds at a never-corrupted leader, as the property states it.}
\deviation{D29}{ghost outputs carried down as auxiliary variables}
\departure{Ghost outputs, the round's bound bit and a gather instance's common core, which Transition Systems~2 and~4 announce on their return labels already, carried down to the implementations as auxiliary variables in the sense of~\cite{AL91}.
  A network component holds each of them and no program reads one.
  A gather instance's core is the field $core$ of the gather network, beside the per-sender sent sets and the corrupted set; the bound bit of a round of the direct implementation is a field of that round's network, and the bound bit of a gather-based round is the whole state of that round's own network, which exchanges no messages; an implementation holds both in the ghost of D30.}
\reason{Announcing the value is what puts binding, a property of the extensions of an execution, on a single trace.
  A refinement carries properties of single traces and no others.}
\deviation{D30}{a parametric per-round ghost held by the network}
\departure{An implementation is parametric in a per-round ghost, held by the network of \leandecl{PLTS.ABA.Implementation.system} alongside the message sets and the corrupted set, and read by no program.
  Two parameters carry it: an update applied on every transition, which replaces the ghost of the round the label names and leaves the other rounds alone, and an output relation on the announced bit guarding the two graded-agreement return transitions, whose three instantiations are Definition~\ref{def:aba-ghost-free}.
  The composed system of the same protocol holds those values in the network components of the round, and the projection relating the two systems reads the ghost onto them.}
\reason{Both implementations instantiate one parametric system, so the ghost is a parameter beside the round transitions; it sits at the network because the value a graded return announces is determined by the round's messages and the corrupted set, which the network holds; and a ghost no program reads, whose output admits a bit at every state, is what the removal of Definition~\ref{def:aba-ghost-free} can delete.}
\deviation{D31}{the coin's threshold counting callers alone}
\departure{The threshold of the resolution \leandeclnamed{resolve}{PLTS.ABA.WCC.Step.resolve} counts callers alone, corrupted callers included, where the source counts the corrupted set beside them.}
\reason{A corrupted process reaches the instance as a caller through its Byzantine call, and a corruption is no access.
  Definition~2.1 of~\cite{ABDY22} defines unpredictability by accesses: the output is unpredictable until at least $d+1$ parties have accessed the coin.
  The count of callers is the count of that definition at $d = f$.
  With $F$ added, the same threshold is reached at the first correct call once $|F| = f$.}
\deviation{D32}{each BIND payload sent by reliable broadcast}
\departure{Each $\mathrm{BIND}$ payload is sent by reliable broadcast, through one bind-broadcast instance per process and gather, where Algorithm~5 of~\cite{AFW25} and Algorithm~4 of~\cite{LeslieBP} send it to all.}
\reason{The adversary of~\cite{AFW25} chooses the corrupted set initially, so a payload sent to all is fixed by its sender's correctness, where under D1 a sender corrupted after its send may equivocate, and the witness that holds the later returns to the core counts committed payloads, which stay fixed whatever happens to their sender.}
\deviation{D33}{both gathers of the round in the binding form}
\departure{Both gathers of the round are the binding form, where the second call of Algorithm~4 of~\cite{AFW25} is non-binding.}
\reason{The round's counting reads the grade off a common core of the second gather, and Remark~22 of~\cite{AFW25} makes the core of a non-binding gather a prophecy variable in the sense of~\cite{AL91}, determined by the completed execution alone, which a forward simulation cannot write at the specification's transition, since the value chosen there must hold in every extension.
  The binding form makes it a history variable, determined by the prefix, and a binding gather is a non-binding gather with one phase more, so every execution of the round is an execution of Algorithm~4 of~\cite{AFW25}.}
\deviation{D34}{Bracha's rounds as published, three echo triggers and the $(n+f)/2$ quorum}
\departure{Bracha's rounds follow Algorithm~1 of~\cite{AFW25}, where Algorithm~6 of the source blueprint fires the echo step on a received $\mathrm{INIT}$ alone and puts $n-f$ at every quorum.
  The echo step fires on a received $\mathrm{INIT}$ from the leader, on received $\mathrm{ECHO}\,m$ messages from more than $(n+f)/2$ senders, or on $f{+}1$ received $\mathrm{VOTE}\,m$ messages, which are the three cases of the guard of \leandeclnamed{echo}{PLTS.ABA.BRB.BrachaAlgorithm.echo}.
  The vote step fires on that $\mathrm{ECHO}$ quorum or on $f{+}1$ received $\mathrm{VOTE}\,m$ messages, and the return takes $2f{+}1$ received $\mathrm{VOTE}\,m$ messages.
  The quorum is \leandecl{PLTS.ABA.Parameters.receivedEchoQuorum}, which is $(n+f)/2 + 1$, and the third message level is named $\mathrm{VOTE}$ here and in the source blueprint where both~\cite{ABDY22} and~\cite{AFW25} write $\mathrm{READY}$.
  The witness the refinement carries is a quorum of received $\mathrm{ECHO}$ messages at that size, and a correct echo of $m$ carries the leader's input unless the leader is corrupted.}
\reason{Algorithm~1 of~\cite{AFW25} is Bracha's algorithm as published; its $\mathrm{ECHO}$ quorum of more than $(n+f)/2$ senders is the witness the refinement carries, two such quorums meeting in a correct sender of the write-once echo, and the three triggers of the echo step are what Bracha's totality rests on.}
\deviation{D35}{the gather transitions carrying the order of the sends}
\departure{The gather transitions carry the order in which Algorithm~5 of~\cite{AFW25} sends, where Algorithm~4 of the source blueprint states each send as a transition firing on its received messages alone.
  The vote requires the sender's own $\mathrm{ECHO}$, the bind call requires the sender's own $\mathrm{VOTE}$ and no earlier bind call, and the return requires the returner's own bind call, which the write-once field $\mathit{sentBind}$ of the gather variables holds.
  The $\mathrm{ECHO}$ payload is the set of entries the input broadcasts have returned to the sender, which is what Algorithm~5 of~\cite{AFW25} and Algorithm~4 of the source blueprint both send.
  The transition fires once that set has $n-f$ entries, which is the source blueprint's guard.}
\reason{The order is that of Algorithm~5's main thread, and each transition's guard on the sender's own send at the level below is the reachability of the wait-until block written as a guard, so the transitions encode the program rather than its handlers alone.}
\deviation{D36}{the input-enabledness loop guarded on its variables}
\departure{The input-enabledness loop of the ABA interface is guarded on the variables it stands beside, \leandeclnamed{inputLoop}{PLTS.ABA.Implementation.ProgramStep.inputLoop} absorbing a call at a process holding an input and the commit transition carrying the label at a process holding none, and Transition System~1 splitting the label the same way, \leandeclnamed{callLoop}{PLTS.ABA.SpecStep.callLoop} looping at a filled ghost entry and \leandeclnamed{callSet}{PLTS.ABA.SpecStep.callSet} writing an empty one, where the source's loop is enabled at every state beside its recording transition, so that a first call may be dropped there and a later one recorded.}
\reason{The call label is enabled at every state in both systems, and every write to the ghost is a first write, so the ghost holds each never-corrupted process's first genuine call.
  That call is the witness Validity names: \leandecl{PLTS.ABA.ValidityTrace} asks of every return of $b$ by a never-corrupted process an earlier $\mathrm{callABA}(id',b)$ at a never-corrupted $id'$, first among $id'$'s $\mathrm{callABA}$ labels in the trace.}
\end{deviations}
The labels D2, D3, D6, D7, D10, D20, D25 and D28 are not in use, and D12 is superseded by D12$'$.


\section{The specifications and their safety}

A specification locates the adversary differently from the implementation it stands for.
An implementation puts processes that read only their own replacement flag beside an explicit network-adversary PLTS, which owns the corrupted set together with its budget and authorises every Byzantine action by the guard $k \in F$ of a component; a specification carries no network component at all, and its adversarial transitions are bare nondeterministic choices resolved by the scheduler.

\subsection{Parameters and the shared alphabet}

The parameters and the label type are fixed first: a refinement is an inclusion of trace-distribution sets, and a trace is a sequence of external labels.
See \leanmodule{Leslie2Protocols/ABA/Vocabulary/Parameters} and \leanmodule{Leslie2Protocols/ABA/Vocabulary/Labels}.

\begin{definition}[Parameters and the coin]
  \label{def:aba-parameters}
  \lean{PLTS.ABA.Parameters}\lean{PLTS.ABA.Parameters.wccPMF}\leanok
  The parameters of the case study are $n$ processes, at most $f$ corrupted with $3f < n$, coin goodness $\varepsilon$ and coin failure mass $\failmass$ subject to $2\varepsilon + \failmass \le 1$, together with the one coin distribution they induce, which both coin resolutions of the development are read off, that of Transition System~1 through the map that drops the bit (D21): \leandecl{PLTS.ABA.Parameters} and \leandecl{PLTS.ABA.Parameters.wccPMF}.
\end{definition}

\begin{definition}[Shared alphabet]
  \label{def:aba-labels}
  \lean{PLTS.ABA.Label}\lean{PLTS.ABA.Label.hiddenAPI}\leanok
  \uses{def:aba-parameters}
  One label type carries the ABA interface, the round-tagged GBCA and WCC interfaces, the corruption broadcast and $\tau$, and the two sub-protocol interfaces make up the hidden alphabet of every composed system: \leandecl{PLTS.ABA.Label} and \leandecl{PLTS.ABA.Label.hiddenAPI}.
  A graded-agreement return $\mathrm{retG}(r,id,out,bnd)$ names the round, the process answered, the graded outcome it receives and the round's bound bit; the outcome is program data and the bit is a ghost output, which no program reads (D29).
\end{definition}


\subsection{The ABA specification and its safety}

Transition System~1 is eight transitions over the state \leandecl{PLTS.ABA.SpecState}, and the flip is its only probabilistic transition.
Safety is proved by induction over those eight transitions under two invariants, one on the state and one reading the label history.
The decision value $val$ is written once, by one transition guarded on $val = \bot$, and the $f+1$ $F$-blind support of the bit it writes is that same transition's guard, so the attribution of the decision to a caller needs no separate invariant conjunct.
The budget pigeonhole then leaves one supporter never corrupted and ghost-recorded as having called the value, which is the caller Validity names.
Both trace predicates quantify over never-corrupted returners: a corrupted process returns an arbitrary bit at an arbitrary time (D23), which no property of the system constrains.
See \leanmodule{Leslie2Protocols/ABA/Specifications/ABA} and \leanmodule{Leslie2Protocols/ABA/Specifications/ABASafety}.

\begin{definition}[The ABA specification, repaired]
  \label{def:aba-spec}
  \lean{PLTS.ABA.SpecStep}\lean{PLTS.ABA.spec}\leanok
  \uses{def:aba-labels}
  Transition System 1 is the repaired ABA specification, eight transitions over the ABA interface on a state carrying a ghost record of the genuine calls, the processes that have returned, the corrupted set, the decision value and a control mode $mode \in \{flipEnabled, decisionEnabled, noTransitionEnabled\}$ (D1, D13, D17, D21, D23), among them the corrupted interface \leandeclnamed{callByzantine}{PLTS.ABA.SpecStep.callByzantine} and \leandeclnamed{retByzantine}{PLTS.ABA.SpecStep.retByzantine} and the flip \leandeclnamed{coinFlip}{PLTS.ABA.SpecStep.coinFlip} under its mixedness guard: \leandecl{PLTS.ABA.SpecStep} and \leandecl{PLTS.ABA.spec}.
\end{definition}


\blueprintfig{fig-ts1-transitions}{Transition System 1, the ABA specification \leandecl{PLTS.ABA.spec}: the four control states, keyed on the mode and the decision value, and the flip branch of \leandeclnamed{coinFlip}{PLTS.ABA.SpecStep.coinFlip}, whose $\tau$-step is drawn with an undirected stem and one edge per outcome — the decision enabled, the flip left enabled, and failure to deliver, naming no coin bit.
The mode $noTransitionEnabled$ enables no silent transition.
Transitions 1 and 2 are drawn on $flipEnabled$ alone.
Transition 1 is enabled at a process whose ghost entry is empty and transition 2 at a process whose entry is filled, so the call label is enabled in every state (D36).
Three of the eight transitions are not drawn: the guarded corruption transition \leandeclnamed{fail}{PLTS.ABA.SpecStep.fail} (D1) and the two corrupted interface transitions \leandeclnamed{callByzantine}{PLTS.ABA.SpecStep.callByzantine} and \leandeclnamed{retByzantine}{PLTS.ABA.SpecStep.retByzantine} (D23), each of which is enabled in every state at a corrupted process and none of which changes the mode.}

\input{nodes-min/thm-aba-spec-safe}

\subsection{The sub-protocol specifications}

Transition Systems~2 and~3 are the specifications of the two sub-protocols, given as round-indexed instances together with the $\mathbb{N}$-indexed families the round loop calls.
The trace form of binding, and its transfer to both implementations, is Theorem~\ref{thm:aba-gbca-binding-trace}.
Both graded-agreement implementations meet Transition System~2 read over the alphabet of a round, along the projection that identifies the Byzantine call and return labels with the specification labels they stand for.
Binding holds of that reading, so every inclusion into it carries binding.
See \leanmodule{Leslie2Protocols/ABA/GBCA/Specification}, \leanmodule{Leslie2Protocols/ABA/GBCA/SpecificationSafety}, \leanmodule{Leslie2Protocols/ABA/GBCA/SpecificationOverRoundAlphabet} and \leanmodule{Leslie2Protocols/ABA/GBCA/BindingOverRoundAlphabet}.

\begin{definition}[The graded-agreement specification]
  \label{def:aba-gbca-spec}
  \lean{PLTS.ABA.GBCA.Step}\lean{PLTS.ABA.GBCA.specFamily}\leanok
  \uses{def:aba-labels, def:idle-family}
  Transition System 2 as a round-indexed instance and as the $\mathbb{N}$-indexed family over it, binding by exclusion under $f{+}1$ support guards (D14, D15, D19): \leandecl{PLTS.ABA.GBCA.specInst} and \leandecl{PLTS.ABA.GBCA.specFamily}.
  Each of the three returns takes a bit $bnd$ under the one guard $1-bnd \in excluded$ and announces it on its label, the $C$-return included, so the round's surviving bit is carried by every return of the round (D19, D29).
\end{definition}


\blueprintfig{fig-gbca-lifecycle}{The GBCA specification instance \leandecl{PLTS.ABA.GBCA.specInst} (Transition System 2): the reachable states quotiented to the pair $(|excluded|, grade)$, the single exclusion \leandeclnamed{bindUnset}{PLTS.ABA.GBCA.Step.bindUnset} under an $f{+}1$ support guard, and the mutually exclusive grade-2 / grade-0 guard.
The seven bit-labelled reachable states are drawn as their four-state quotient by the exchange of $0$ and $1$; $(\emptyset,2)$ and $(\emptyset,0)$ are unreachable, every graded return requiring $excluded \neq \emptyset$.
The graded returns are \leandeclnamed{retGrade2}{PLTS.ABA.GBCA.Step.retGrade2}, \leandeclnamed{retGrade1}{PLTS.ABA.GBCA.Step.retGrade1} and \leandeclnamed{retGrade0}{PLTS.ABA.GBCA.Step.retGrade0}, and every $\mathrm{InputSupport}$ count in their guards is $f{+}1$ distinct callers-or-corrupted (D15).
Each of the three carries the round's bound bit on its label, drawn beside it as $bnd = 1{-}b$: a return announces the complement of the excluded bit, so every return of a round announces one bit (D29).}

\begin{theorem}[Binding and graded agreement of the GBCA specification]
  \label{thm:gbca-spec-binding}
  \lean{PLTS.ABA.GBCA.excluded_card_le_one}\lean{PLTS.ABA.GBCA.retG_value_agree}\lean{PLTS.ABA.GBCA.retG_grade_exclusive}\lean{PLTS.ABA.GBCA.retG_bound_agree}\lean{PLTS.ABA.GBCA.retG_value_eq_bound}\lean{PLTS.ABA.GBCA.retGrade0_excluded_nonempty}\leanok
  \uses{def:aba-gbca-spec}
  Along one execution of the round-$r$ GBCA specification instance, any two returns that hand out a bit hand out the same bit, any two returns announce the same bit, and a return that hands out a bit announces the bit it hands out: \leandecl{PLTS.ABA.GBCA.retG_value_agree}, \leandecl{PLTS.ABA.GBCA.retG_bound_agree} and \leandecl{PLTS.ABA.GBCA.retG_value_eq_bound}.
  No execution carries both a grade-2 return and a grade-0 return, which is the second clause of graded agreement: \leandecl{PLTS.ABA.GBCA.retG_grade_exclusive}, read off the trace as \leandecl{PLTS.ABA.GBCA.specInst_grade_agree}.
  The exclusion set holds at most one bit at every state, and a grade-0 return leaves it nonempty at every later state, which is the Graded Binding witness: \leandecl{PLTS.ABA.GBCA.excluded_card_le_one} and \leandecl{PLTS.ABA.GBCA.retGrade0_excluded_nonempty}.
\end{theorem}

\begin{proof}\leanok
  \uses{def:aba-gbca-spec}
  Monotonicity of the exclusion set and the bound $|excluded| \le 1$ give every clause but the last, and monotonicity of the write-once grade gives that one; see \leandecl{PLTS.ABA.GBCA.retG_value_agree} and \leandecl{PLTS.ABA.GBCA.retG_grade_exclusive}.
\end{proof}

\begin{theorem}[Validity of the GBCA specification, safety half]
  \label{thm:gbca-spec-validity}
  \lean{PLTS.ABA.GBCA.specInst_validity}\leanok
  \uses{def:aba-gbca-spec, thm:trace-support}
  On a trace of the GBCA specification instance whose every round-$r$ call carries $v$ unless its caller is corrupted somewhere along the trace, every round-$r$ return hands out $v$: \leandecl{PLTS.ABA.GBCA.specInst_validity}.
  Neither a $C$-return nor a $B$-return occurs on such a trace, so every return is an $A$-return of $v$: \leandecl{PLTS.ABA.GBCA.specInst_no_retC} and \leandecl{PLTS.ABA.GBCA.specInst_no_retB}.
\end{theorem}

\begin{proof}\leanok
  \uses{def:aba-gbca-spec, thm:trace-support}
  See the Lean proof of \leandecl{PLTS.ABA.GBCA.specInst_validity}.
\end{proof}


The WCC instance, over \leandecl{PLTS.ABA.WCC.SpecState}, resolves its coin once into four outcomes, by a silent transition enabled once the callers number more than $f$.
That resolution is the second and last non-Dirac transition of the development.
Its value domain is \leandecl{PLTS.ABA.CoinValue}, ranging over $\{0,1,\bot,\top,\dagger\}$, where $\bot$ is unresolved, $\top$ is an adversarially resolved outcome and $\dagger$ is a resolution that delivers to no process (D17).
The resolution is the one silent transition of the instance, and the scheduler places it at any state after the count exceeds $f$.
Every other transition is an access, a return or a corruption.
See \leanmodule{Leslie2Protocols/ABA/Specifications/WCC}.

\begin{definition}[The coin specification]
  \label{def:aba-wcc-spec}
  \lean{PLTS.ABA.WCC.Step}\lean{PLTS.ABA.WCC.specFamily}\leanok
  \uses{def:aba-labels, def:aba-parameters, def:idle-family}
  Transition System 3 as a round-indexed instance and as the $\mathbb{N}$-indexed family over it, its call carrying the two transitions \leandeclnamed{call}{PLTS.ABA.WCC.Step.call} and \leandeclnamed{callLoop}{PLTS.ABA.WCC.Step.callLoop}: \leandecl{PLTS.ABA.WCC.specInst} and \leandecl{PLTS.ABA.WCC.specFamily}.
  The resolution \leandeclnamed{resolve}{PLTS.ABA.WCC.Step.resolve} is a silent transition, enabled at $val = \bot$ once the count \leandeclnamed{threshold}{PLTS.ABA.WCC.SpecState.threshold}, taken over the callers alone (D31), exceeds $f$, and it draws $val$ by the distribution of Definition~\ref{def:aba-parameters}.
  The loop \leandeclnamed{callLoop}{PLTS.ABA.WCC.Step.callLoop} records no caller, and the scheduler places the resolution at any state after the count exceeds $f$.
\end{definition}


\blueprintfig{fig-wcc}{The WCC specification instance \leandecl{PLTS.ABA.WCC.specInst} (Transition System 3): the four-outcome coin resolution and the return of the common bit or an adversary-chosen bit.
The resolution \leandeclnamed{resolve}{PLTS.ABA.WCC.Step.resolve} is a silent transition, enabled at $val = \bot$ once the callers number more than $f$, and every call is \leandeclnamed{call}{PLTS.ABA.WCC.Step.call} or \leandeclnamed{callLoop}{PLTS.ABA.WCC.Step.callLoop}.
The branch $val{=}\dagger$ of mass $\failmass$ enables no return (D17), \leandeclnamed{ret}{PLTS.ABA.WCC.Step.ret} having a positive guard.}

\section{The protocol}
\label{sec:aba-protocol}

What runs is $n$ programs, a network holding the message sets and the corrupted set with its budget, and the common coin held at specification level.
A program reads its own replacement flag and nothing else about corruption: not the corrupted set, not the budget, not another process's status (D23).
The synchronisations the shared alphabet cannot name are carried by an extended alphabet, hidden to $\tau$ and then discharged by the restriction, so what leaves the outermost hiding is the alphabet of the ABA specification.
The adversary holds one thing more, per round, which no program reads: the ghost of D30, written on every transition by a parameter of the implementation and read out by a second parameter that guards the two graded-agreement return transitions.
That is where the bound bit a return announces comes from, and it is the second thing a graded-agreement implementation supplies.
A system below the meeting point therefore belongs to one implementation or the other, and carries a superscript naming its source: $\abdy{\mathrm{protocol}}(P)$ and $\abdy{\mathrm{composed}}(P)$ after~\cite{ABDY22}, and $\afw{\mathrm{protocol}}(P)$ and $\afw{\mathrm{composed}}(P)$ after~\cite{AFW25}.
The systems the two share, $\mathrm{hybrid}(P)$ and $\mathrm{spec}(P)$, carry none.

The ghost is written by every transition and read by the two graded-agreement return transitions alone, so the one thing it can affect is the bit such a return announces.
That label belongs to the hidden alphabet of Definition~\ref{def:aba-labels}, and the outermost hiding removes it, so the bit reaches no trace of the protocol.
Definition~\ref{def:aba-ghost-free} is the same system with the ghost deleted and the announcement left to the adversary, and Theorem~\ref{thm:aba-ghost-removal} is the equality of achievable trace distributions between the two, from which every conclusion of either chain holds of the ghost-free system.

The control state beneath the implementation holds one process's round loop, and no field of it names another process or holds the corrupted set.
See \leanmodule{Leslie2Protocols/ABA/Implementation/Alphabet}, \leanmodule{Leslie2Protocols/ABA/Implementation/System}, \leanmodule{Leslie2Protocols/ABA/Implementation/StepCases}, \leanmodule{Leslie2Protocols/ABA/Implementation/NetworkStateWritesAndRemovals}, \leanmodule{Leslie2Protocols/ABA/Implementation/CompositeTransitions} and \leanmodule{Leslie2Protocols/ABA/Vocabulary/RoundLoop}.

\begin{definition}[The state one process carries between rounds]
  \label{def:aba-round-loop-state}
  \lean{PLTS.ABA.RoundLoopState}\leanok
  The control state of one process in the round loop holds its input bit, its estimate, its round counter, its phase in the round's calls and returns, the round's graded outcome and a returned flag, and names no other process: \leandecl{PLTS.ABA.RoundLoopState}.
\end{definition}

\begin{definition}[The implementation of a protocol]
  \label{def:aba-implementation}
  \lean{PLTS.ABA.Implementation.NetworkEvent}\lean{PLTS.ABA.Implementation.actsAt}\lean{PLTS.ABA.Implementation.coinLabelMap}\lean{PLTS.ABA.Implementation.RoundVariablesMap}\lean{PLTS.ABA.Implementation.IsRoundVariables}\lean{PLTS.ABA.Implementation.NetworkState}\lean{PLTS.ABA.Implementation.roundOf}\lean{PLTS.ABA.Implementation.NetworkState.writeGhost}\lean{PLTS.ABA.Implementation.roundOwn}\lean{PLTS.ABA.Implementation.ProgramStep}\lean{PLTS.ABA.Implementation.NetworkStep}\lean{PLTS.ABA.Implementation.IsRoundStep}\lean{PLTS.ABA.Implementation.system}\leanok
  \uses{def:aba-labels, def:aba-parameters, def:aba-round-loop-state, def:relabel, def:synchronised-product, def:process-algebra, def:idle-family}
  A protocol as it runs: $n$ programs, one per process, beside one network holding the message sets, the DECIDED sets and the corrupted set with its budget, beside the common coin.
  The system is parametric in the graded-agreement implementation --- its message type, its round's own events, its per-process per-round variables and its transitions --- and in a per-round ghost the adversary holds, with an update \leandecl{PLTS.ABA.Implementation.NetworkState.writeGhost} applied on every transition and an output relation on the announced bit guarding the two graded-agreement return transitions (D30), so the round loop, the DECIDED sets, the coin's call and return, corruption, the adversary and the composition pipeline are written once: \leandecl{PLTS.ABA.Implementation.system}.
  What an implementation supplies about its own transitions, and what the lemmas that read a transition off its label consume, is \leandecl{PLTS.ABA.Implementation.IsRoundStep}.
\end{definition}


\section{The ABDY implementation}
\label{sec:aba-abdy}

The first of the two implementations of graded agreement, and the protocol built on it.

\subsection{The graded-agreement round}

The first implementation of graded agreement calls no sub-protocol: it is ABDY22's Algorithm~6 transcribed in full, over five message levels.
The depth of those levels is what the grade-$1$ witness buys (D18): $f{+}1$ received $\mathrm{BIND}\,v$ messages exceed the corruption budget, so a correct sender stands behind every grade-$1$ output, and that sender's own wait condition is the $n{-}f$ $\mathrm{VOTE}\,v$ quorum the binding argument of~\cite{ABDY22} counts.
The round variables hold one process's state in the graded-agreement round, and the definition below gives the algorithm that round runs, which is the form Theorem~\ref{thm:aba-gbca-composition-by-abdy-refines-specification} is proved against.
Its state is the pair of the $n$ processes' round variables and the round's own network state, and its guards read that pair through four accessors.
Beside the message sets and the corrupted set that network state holds the round's bound bit, written at the round's first return and announced on every return of the round (D29).
It is computed from the round's sent sets, the corrupted set and the graded outcome, and it enters no guard of the algorithm.
Definitions~\ref{def:aba-programs} and~\ref{def:aba-network} below are where each of those guards is written against one process's variables.
The round variables and the round's network state are in \leanmodule{Leslie2Protocols/ABA/GBCA/ABDY/MessagesAndVariables}, and the algorithm is in \leanmodule{Leslie2Protocols/ABA/GBCA/ABDY/Algorithm}.

\begin{definition}[Round messages and variables]
  \label{def:aba-round-messages}
  \lean{PLTS.ABA.GBCA.ByABDY.Message}\lean{PLTS.ABA.GBCA.ByABDY.ProcessVariables}\leanok
  The messages of one graded-agreement round run over five message levels, and the round state of one process holds the input its call records, its write-once send fields and its return flag, no field of which names another process: \leandecl{PLTS.ABA.GBCA.ByABDY.Message} and \leandecl{PLTS.ABA.GBCA.ByABDY.ProcessVariables}.
\end{definition}

\begin{definition}[The algorithm of the round's graded-agreement composition]
  \label{def:aba-gbca-algorithm-by-abdy}
  \lean{PLTS.ABA.GBCA.ByABDY.boundOf}\lean{PLTS.ABA.GBCA.ByABDY.Algorithm}\lean{PLTS.ABA.GBCA.ByABDY.composition_projects}\leanok
  \uses{def:aba-labels, def:aba-parameters, def:aba-round-messages}
  The algorithm one graded-agreement round runs, as a relation on the pair of local states that hold a round's state --- the $n$ processes' round variables beside the round's network state --- transcribing ABDY22's Algorithm~6 over the five message levels, one transition per line (D18): \leandecl{PLTS.ABA.GBCA.ByABDY.Algorithm}.
  At a label of the round's interface, every transition of the round's composition is the algorithm's transition at the specification label that interface label projects to, one step for one step: \leandecl{PLTS.ABA.GBCA.ByABDY.composition_projects}.
  The network state holds one field more, the round's bound bit: the three return transitions announce \leandecl{PLTS.ABA.GBCA.ByABDY.boundOf} of the round's sent sets, the corrupted set and the outcome where the round holds no bit, and the bit it holds otherwise, and write it back (D29).
\end{definition}


\blueprintfig{fig-gbca-impl}{The round's graded-agreement composition \leandecl{PLTS.ABA.GBCA.ByABDY.composition}, running ABDY22's Algorithm~6 (D18): the per-process round automaton and, framed, the authenticated set-based network of D5.
Each return box carries the bit its label announces: the grade-2 and grade-1 returns the bit they hand out, the grade-0 return the bit read off the sent sets.
The network holds that bit for the round as a ghost, written at the first return and announced at every return, and no program's variables hold it (D29).}

\subsection{The programs and the protocol}

The protocol of Section~\ref{sec:aba-protocol} at that graded-agreement implementation is the $n$ programs, each carrying its variables in every round it has touched, beside the one adversary holding what no process may see.

Corruption is the one synchronised step $\mathrm{fail}(j)$, on a visible label.
Its halves are the network's guarded insertion \leandeclnamed{fail}{PLTS.ABA.Implementation.NetworkStep.fail}, the write of the process's own flag $corrupted$ by \leandeclnamed{failSelf}{PLTS.ABA.Implementation.ProgramStep.failSelf}, every other program idle, and each instance of the common coin inserting $j$ into its own copy of the corrupted set.
Every correct transition reads the flag, and the replaced program \leandeclnamed{corruptedIdle}{PLTS.ABA.Implementation.ProgramStep.corruptedIdle} self-loops on every label other than $\tau$ and those of \leandeclnamed{actsAt}{PLTS.ABA.Implementation.actsAt}\,$j$, on which it has no transition, so a corrupted process neither calls into nor returns from a round at the implementation, while the ABA-level corrupted interface, the network's Byzantine return under $j \in F$, fires.
The clause \leandeclnamed{corrupted_F}{PLTS.ABA.Invariant.corrupted_F} of the invariant of Theorem~\ref{thm:aba-hybrid-spec} ties the flag to the set.
No transition emits $\mathrm{fail}$: the scheduler fires it wherever it is enabled, and the trace predicates read it off the trace.

\begin{definition}[The program of one process]
  \label{def:aba-programs}
  \lean{PLTS.ABA.GBCA.ByABDY.RoundRecord}\lean{PLTS.ABA.RoundLoopRecord}\lean{PLTS.ABA.ABDY.RoundRecordMap}\lean{PLTS.ABA.ABDY.ProcessRecord}\lean{PLTS.ABA.Composition.NetworkEvent}\lean{PLTS.ABA.Composition.ExtendedLabel}\lean{PLTS.ABA.Implementation.networkEventLabels}\lean{PLTS.ABA.ABDY.RoundStep}\lean{PLTS.ABA.ABDY.ABAProgramStep}\lean{PLTS.ABA.ABDY.ABAProgram}\leanok
  \uses{def:aba-labels, def:aba-round-messages, def:aba-round-loop-state, def:aba-implementation, def:relabel, def:synchronised-product}
  The program of process $j$ of the protocol, Definition~\ref{def:aba-implementation} at ABDY22's implementation of Definition~\ref{def:aba-round-messages}, on the state $\mathrm{ProcessRecord}(n) = \mathrm{RoundLoopRecord}(n) \times \mathrm{RoundRecordMap}(n)$ --- the round-loop record beside the round record of every round the process has touched and its termination flag, which the round advance does not reset and which the transition $terminate$ leaves standing (D22): \leandecl{PLTS.ABA.ABDY.ABAProgram}.
  The round's own transitions are the fourteen of \leandecl{PLTS.ABA.ABDY.RoundStep}: the eight multicasts of those levels, the delivery, the two calls and the three returns.
  It reads no copy of the corrupted set and no budget, only the replacement flag \leandecl{PLTS.ABA.RoundLoopRecord.corrupted} of its own record, which \leandeclnamed{failSelf}{PLTS.ABA.Implementation.ProgramStep.failSelf} writes and every other transition requires to be down; in place of those transitions the replaced program is the self-loop \leandeclnamed{corruptedIdle}{PLTS.ABA.Implementation.ProgramStep.corruptedIdle} (D23).
\end{definition}

\begin{definition}[The network]
  \label{def:aba-network}
  \lean{PLTS.ABA.ABDY.NetworkState}\lean{PLTS.ABA.ABDY.NetworkState.corrupt}\lean{PLTS.ABA.ABDY.abdyGhostStep}\lean{PLTS.ABA.ABDY.abdyGhostOutput}\lean{PLTS.ABA.ABDY.NetworkStep}\lean{PLTS.ABA.ABDY.network}\leanok
  \uses{def:aba-programs, def:aba-labels, def:aba-parameters}
  The one local state holding what no process may see: the per-sender sent sets of the round messages and of the DECIDED payloads, and the corrupted set with its budget, which guards its $\mathrm{fail}$ transition (D1) and authorises the five Byzantine call and return transitions (D11), the two injections and the Byzantine return $\mathrm{retByzantine}$ (D23): \leandecl{PLTS.ABA.ABDY.network}.
  Its ghost for a round is one bit, the round's bound bit: a return records the bit its own label announces where the round holds none, and the two return transitions fire only with that bit equal to \leandecl{PLTS.ABA.ABDY.abdyGhostOutput}, which is the bit the round holds or \leandecl{PLTS.ABA.GBCA.ByABDY.boundOf} of the round's messages (D29, D30).
\end{definition}

\begin{definition}[The protocol]
  \label{def:aba-protocol}
  \lean{PLTS.System.mapIdle}\lean{PLTS.ABA.Implementation.coinLabelMap}\lean{PLTS.ABA.Composition.coinOverRoundAlphabet}\lean{PLTS.ABA.ABDY.protocolExtended}\lean{PLTS.ABA.ABDY.protocolHidden}\lean{PLTS.ABA.ABDY.protocol}\leanok
  \uses{def:aba-network, def:aba-programs, def:aba-implementation, def:aba-wcc-spec, def:aba-labels, def:relabel, def:synchronised-product, def:process-algebra}
  The pipeline of Definition~\ref{def:aba-implementation} at ABDY22's implementation: the $n$ programs beside the network and the common coin, the coin read over the extended alphabet along a partial label map, under the two hidings: \leandecl{PLTS.ABA.ABDY.protocol}.
\end{definition}


\blueprintfig{fig-protocol}{The protocol \leandecl{PLTS.ABA.ABDY.protocol}: the $n$ programs \leandecl{PLTS.ABA.ABDY.ABAProgram} under $\mathrm{synchronisedProduct}$, the network \leandecl{PLTS.ABA.ABDY.network} over the state \leandecl{PLTS.ABA.ABDY.NetworkState}, and the common coin beside the group.
The adversary also holds the ghost $ghost(r)$ of each round, that round's bound bit, which no program reads (D30).
The dashed vertical is Theorem~\ref{thm:aba-protocol-composed}, along the Dirac lift of $\mathrm{ProtocolRelation}$.}

\subsection{The composed system}

A sub-protocol may be exchanged for its specification where it is a component of a composition, and the protocol does not present graded agreement as one.
The component boundary is where the chain passes from implementation to specification: $\abdy{\mathrm{protocol}}(P)$ is the system that runs, everything above it is specification, and the first step is accordingly an inclusion.
The composed system holds one graded-agreement instance per round at every moment, and a round instance carries no termination flag, so it answers a send or a delivery at a process the protocol has terminated (D22).
Graded agreement is a parallel component of the composed system, which is what the substitution congruence consumes.
See \leanmodule{Leslie2Protocols/ABA/GBCA/ABDY/Components}, \leanmodule{Leslie2Protocols/ABA/GBCA/ABDY/Composition}, \leanmodule{Leslie2Protocols/ABA/GBCA/ABDY/CompositionStepCases}, \leanmodule{Leslie2Protocols/ABA/GBCA/ABDY/Algorithm}, \leanmodule{Leslie2Protocols/ABA/GBCA/ABDY/RoundFamilyOwnedLabels}, \leanmodule{Leslie2Protocols/ABA/Composition/Hybrid} and \leanmodule{Leslie2Protocols/ABA/ABDY/Composition}.

\begin{definition}[The round's graded-agreement instance]
  \label{def:aba-gbca-instance-by-abdy}
  \lean{PLTS.ABA.GBCA.ByABDY.GBCAEvent}\lean{PLTS.ABA.GBCA.ByABDY.GBCALabel}\lean{PLTS.ABA.GBCA.ByABDY.gbcaEvents}\lean{PLTS.ABA.GBCA.ByABDY.GBCAProgramStep}\lean{PLTS.ABA.GBCA.ByABDY.gbcaProgram}\lean{PLTS.ABA.GBCA.ByABDY.NetworkState}\lean{PLTS.ABA.GBCA.ByABDY.NetworkState.corrupt}\lean{PLTS.ABA.GBCA.ByABDY.GBCANetworkStep}\lean{PLTS.ABA.GBCA.ByABDY.GBCANetwork}\lean{PLTS.ABA.GBCA.ByABDY.RoundState}\lean{PLTS.ABA.GBCA.ByABDY.compositionExtended}\lean{PLTS.ABA.GBCA.ByABDY.composition}\lean{PLTS.ABA.GBCA.ByABDY.roundOwnsLabel}\lean{PLTS.ABA.GBCA.ByABDY.corruptionAct}\lean{PLTS.ABA.GBCA.ByABDY.gbcaInstanceFamily}\leanok
  \uses{def:aba-programs, def:aba-network, def:aba-round-messages, def:relabel, def:synchronised-product, def:idle-family, def:process-algebra}
  One round of graded agreement cut into $n$ local graded-agreement programs beside that round's own network, together with the $\mathbb{N}$-indexed family of such rounds that is the graded-agreement component of the protocol: \leandecl{PLTS.ABA.GBCA.ByABDY.composition} and \leandecl{PLTS.ABA.GBCA.ByABDY.gbcaInstanceFamily}.
  The network holds the round's per-sender message sets, the corrupted set and the round's bound bit, the last a ghost output the return rows write and announce (D29).
\end{definition}

\begin{definition}[The composed system]
  \label{def:aba-composed}
  \lean{PLTS.ABA.Composition.RoundLoopStep}\lean{PLTS.ABA.Composition.roundLoopProgram}\lean{PLTS.ABA.Composition.ABANetworkState}\lean{PLTS.ABA.Composition.ABANetworkState.corrupt}\lean{PLTS.ABA.Composition.ABANetworkStep}\lean{PLTS.ABA.Composition.ABANetwork}\lean{PLTS.ABA.Composition.ComposedState}\lean{PLTS.ABA.ABDY.composedExtended}\lean{PLTS.ABA.ABDY.composedHidden}\lean{PLTS.ABA.ABDY.composed}\leanok
  \uses{def:aba-gbca-instance-by-abdy, def:aba-programs, def:aba-network, def:aba-protocol, def:aba-round-loop-state, def:relabel, def:synchronised-product, def:process-algebra}
  The protocol read as a composition of four components --- the family of round instances, the $n$ round loops, the ABA network and the common coin --- with the graded-agreement family in the refinable position and one round instance per round at every moment, answering also at a process the protocol has terminated (D22): \leandecl{PLTS.ABA.ABDY.composed}.
  A round loop carries the replacement flag and the self-loop \leandeclnamed{corruptedIdle}{PLTS.ABA.Composition.RoundLoopStep.corruptedIdle} of the replaced program, the ABA network the guarded $\mathrm{fail}$ transition and the Byzantine return \leandeclnamed{retByzantine}{PLTS.ABA.Composition.ABANetworkStep.retByzantine} (D1, D23).
\end{definition}


\blueprintfig{fig-composed}{The composed system \leandecl{PLTS.ABA.ABDY.composed} and the substitution: the four components of Definition~\ref{def:aba-composed} inside the two hidings of the protocol's pipeline.
The round instance \leandecl{PLTS.ABA.GBCA.ByABDY.composition} is drawn open, the $n$ graded-agreement programs of round $r$ beside that round's own network, with the multicast and the delivery hidden inside the instance.
That network holds the round's per-sender sent sets, the corrupted set and the round's bound bit, a ghost the three return transitions write and announce (D29).
The left position is the refinable one: Theorem~\ref{thm:aba-substitution} replaces the family of round instances by the family of round specifications and leaves the other three components as they stand.}

\subsection{The round loop}

The round loop is the core algorithm of ABA, and it is what calls the other two.
Its external return is guarded by an $n-f$ DECIDED quorum, and that count is what excludes disagreement: two returns of different bits would need two sender sets of size $n-f$, and any two such sets meet in a correct process.
Its transitions are written against one process's own variables in Definition~\ref{def:aba-programs}, and the state below is what those transitions move, read across the two components that hold it.
Both implementations run the same round loop over the same ABA state.
See \leanmodule{Leslie2Protocols/ABA/Vocabulary/RoundLoop} and \leanmodule{Leslie2Protocols/ABA/Composition/ABAState}.

\begin{definition}[The per-process variables beside the ABA network]
  \label{def:aba-state}
  \lean{PLTS.ABA.ABAState}\lean{PLTS.ABA.ABAState.processes}\lean{PLTS.ABA.ABAState.decidedSent}\lean{PLTS.ABA.ABAState.decidedReceived}\lean{PLTS.ABA.ABAState.corrupted}\lean{PLTS.ABA.ABAState.F}\lean{PLTS.ABA.ABAState.decidedCount}\lean{PLTS.ABA.ABAState.setProcessVariables}\lean{PLTS.ABA.ABAState.stepRound}\lean{PLTS.ABA.ABAState.sendDecidedOnGrade2}\lean{PLTS.ABA.ABAState.sendDecided}\lean{PLTS.ABA.ABAState.deliverDecided}\lean{PLTS.ABA.ABAState.corrupt}\leanok
  \uses{def:aba-labels, def:aba-round-loop-state, def:aba-programs, def:aba-composed}
  The $n$ processes' round-loop variables beside the ABA network, read as one object by five accessors that gather the control states, the received DECIDED payloads, the replacement flags (D23), the multicast sent sets and the corrupted set the two components hold apart: \leandecl{PLTS.ABA.ABAState}.
\end{definition}


\blueprintfig{fig-core-phases}{The round loop \leandecl{PLTS.ABA.Composition.roundLoopProgram}: the per-process cycle of six phases of calls and returns, the multicast of $\langle\mathrm{DECIDED},b\rangle$ at $\mathrm{toSendDecided}$ after a grade-2 coin return, and the DECIDED sets guarding $\mathrm{retABA}$.
A corruption leaves the machine at whatever phase it reaches: every edge drawn carries the guard that the program is unreplaced, so the phase field is unchanged from the corruption on (D23).
The \texttt{inputLoop} self-loop absorbs $\mathrm{callABA}$ at a process whose input is set (D36), so it stands at every phase the \texttt{input} edge leads to and is drawn at $\mathrm{toCallG}$ alone.}

\subsection{The refinements}

The work that carries this implementation into Transition System~2 is done once, on one round instance.
The relation reads the specification's call entries, returns and corrupted set straight off the composed state, and witnesses the two fields the protocol does not hold, $excluded$ and $grade$, by \emph{patterns of received messages}: $F$-blind counts on the delivered sets, monotone along every execution, that witness the value of such a field; an exclusion is fired on demand, as a two-step weak run under a return that needs it, which is why the obligation is discharged by a weak transition rather than by a single step.
The exclusion set is also an equation on the round's bound bit: it is empty while the bit is unwritten and the singleton of the bit's complement once a return has written it.
See \leanmodule{Leslie2Protocols/ABA/GBCA/ABDY/Invariant}, \leanmodule{Leslie2Protocols/ABA/GBCA/ABDY/ExclusionWitness}, \leanmodule{Leslie2Protocols/ABA/GBCA/ABDY/SpecificationRelation}, \leanmodule{Leslie2Protocols/ABA/GBCA/ABDY/RefinesSpecification} and \leanmodule{Leslie2Protocols/ABA/GBCA/ABDY/Binding}.

\begin{theorem}[The round's graded-agreement composition refines its specification]
  \label{thm:aba-gbca-composition-by-abdy-refines-specification}
  \lean{PLTS.ABA.GBCA.ByABDY.refinesSpecification}\lean{PLTS.ABA.GBCA.ByABDY.specificationRelation_transition}\lean{PLTS.ABA.GBCA.ByABDY.composition_projects}\lean{PLTS.ABA.GBCA.specificationLabelMap}\lean{PLTS.ABA.GBCA.specificationOverRoundAlphabet}\lean{PLTS.ABA.GBCA.ByABDY.specificationRelation_init}\lean{PLTS.ABA.GBCA.ByABDY.refinesSpecification_failAct}\lean{PLTS.ABA.GBCA.ByABDY.specificationCorruptionAct}\leanok
  \uses{def:aba-gbca-instance-by-abdy, def:aba-gbca-algorithm-by-abdy, def:aba-gbca-spec, def:forwardsim}
  Each round's composition is forward-simulated by that round's GBCA specification instance, read over the round's interface, along a relation that certifies the exclusion set and the grade by receipt patterns, fires an exclusion on demand, holds at the initial states and survives the corruption broadcast: \leandecl{PLTS.ABA.GBCA.ByABDY.refinesSpecification}.
\end{theorem}

\begin{proof}\leanok
  \uses{def:aba-gbca-instance-by-abdy, def:aba-gbca-algorithm-by-abdy, def:aba-gbca-spec, def:forwardsim, def:weakstep}
  See the Lean proof of \leandecl{PLTS.ABA.GBCA.ByABDY.refinesSpecification}, which reads a transition of the composition as one of the algorithm at \leandecl{PLTS.ABA.GBCA.ByABDY.composition_projects} and matches that transition at \leandecl{PLTS.ABA.GBCA.ByABDY.specificationRelation_transition}.
\end{proof}


\blueprintfig{fig-gbca-sim}{The forward simulation \leandecl{PLTS.ABA.GBCA.ByABDY.refinesSpecification} of the round's composition, along the relation \leandecl{PLTS.ABA.GBCA.ByABDY.specificationRelation}: the algorithm's transitions above, specification-instance states below, dashed verticals the relation.
A return whose bit the specification has not yet excluded is matched by a weak run, mapping to the weak step $\mathrm{bindUnset};\mathrm{ret}$ whose excluded bit is supplied by the return's own exclusion witness; a return whose exclusion has already happened maps to a single graded specification return.
The relation also reads $excluded$ off the network state's bound bit: empty while the bit is unwritten, and the singleton of the bit's complement once a return has written it.}

\begin{theorem}[The protocol under its composed system]
  \label{thm:aba-protocol-composed}
  \lean{PLTS.ABA.ABDY.ProtocolState}\lean{PLTS.ABA.ABDY.ProtocolRelation}\lean{PLTS.ABA.ABDY.protocolRelation_mk}\lean{PLTS.ABA.ABDY.protocolRelation_init}\lean{PLTS.ABA.ABDY.protocolSimulation}\lean{PLTS.ABA.ABDY.protocol_composed}\leanok
  \uses{def:aba-composed, def:aba-programs, def:aba-network, def:aba-protocol, def:forwardsim, def:coupling, def:achievable, thm:forward-sound}
  The protocol is forward-simulated by its composed system along the Dirac lift of $\mathrm{ProtocolRelation}(P)$, the relation that fixes every coordinate of a composed state that a state of the protocol determines, in five unguarded conjuncts, and hence $\operatorname{ATD}(\abdy{\mathrm{protocol}}(P)) \subseteq \operatorname{ATD}(\abdy{\mathrm{composed}}(P))$; the inclusion is one-directional because a round instance answers the Byzantine graded-agreement transitions, for which no program has a transition (D11), and answers a send or a delivery at a process the protocol has terminated (D22): \leandecl{PLTS.ABA.ABDY.protocolSimulation}, \leandecl{PLTS.ABA.ABDY.protocol_composed}.
  A corrupted process's replaced program is matched loop for loop, and its corrupted return by the ABA network's Byzantine transition (D23).
\end{theorem}

\begin{proof}\leanok
  \uses{def:aba-composed, def:aba-programs, def:aba-network, def:aba-protocol, def:forwardsim, def:coupling, def:achievable, thm:forward-sound}
  See the Lean proof of \leandecl{PLTS.ABA.ABDY.protocolSimulation}, then soundness (Theorem~\ref{thm:forward-sound}).
\end{proof}


\section{The hybrid and the ABDY chain}

Replacing the graded-agreement family by its specification leaves $\mathrm{hybrid}(P)$, which is where the two implementations meet.

\subsection{The substitution}

The graded-agreement family occupies the refinable position of the composed system, so exchanging it for its specification is one family congruence followed by the precongruences of the pipeline.
What the idealisation removes is the round's network and every Byzantine power over it; the DECIDED sets, the corrupted set with the budget $|F| < f$ as the guard on the $\mathrm{fail}$ transition, and the authorisation of every Byzantine call and return transition are outside the exchanged component and survive unchanged.
The system that results, $\mathrm{hybrid}(P)$, is where the two implementations meet: both reach it, and every step above it is proved once and shared.
It is a system for every type of round message, so each implementation reaches it at its own: ABDY22's five message levels, and the empty type for the gather-based round, which exchanges its messages inside its gather instances.
See \leanmodule{Leslie2Protocols/Framework/FamilySimulation}, \leanmodule{Leslie2Protocols/ABA/Composition/Hybrid}, \leanmodule{Leslie2Protocols/ABA/ABDY/Composition} and \leanmodule{Leslie2Protocols/ABA/ABDY/Substitution}.

\begin{theorem}[Family congruence and the graded-agreement family refinement]
  \label{thm:gbca-family-refines}
  \lean{PLTS.ForwardSimulation.family}\lean{PLTS.ABA.gbcaSpecificationFamily}\lean{PLTS.ABA.substitutionRelationFamily}\lean{PLTS.ABA.familySubstitution}\lean{PLTS.ABA.familySubstitutionSimulation}\leanok
  \uses{thm:aba-gbca-composition-by-abdy-refines-specification, def:aba-gbca-instance-by-abdy, def:idle-family, lem:hyperstep-weakstep}
  Per-instance forward simulations lift to the $\mathbb{N}$-indexed families, and the lift makes the family of round instances forward-simulated by the family of round specifications: \leandecl{PLTS.ForwardSimulation.family}, \leandecl{PLTS.ABA.familySubstitutionSimulation}.
\end{theorem}

\begin{proof}\leanok
  \uses{thm:aba-gbca-composition-by-abdy-refines-specification, def:idle-family, lem:hyperstep-weakstep}
  See the Lean proofs of \leandecl{PLTS.ForwardSimulation.family} and \leandecl{PLTS.ABA.familySubstitutionSimulation}.
\end{proof}

\begin{definition}[The hybrid]
  \label{def:aba-hybrid}
  \lean{PLTS.ABA.HybridState}\lean{PLTS.ABA.hybridExtended}\lean{PLTS.ABA.hybrid}\leanok
  \uses{def:aba-composed, def:aba-gbca-instance-by-abdy, def:aba-gbca-spec, def:idle-family, def:relabel, def:process-algebra}
  The composed system with its graded-agreement component replaced by the family of round specifications, the other three components and the whole pipeline unchanged: \leandecl{PLTS.ABA.hybrid}.
  It is a system $\mathrm{hybrid}(P, M)$ for every type $M$ of the messages a graded-agreement round exchanges: the ABDY chain takes $\mathrm{GBCA.ByABDY.Message}$ for it, the gather-based chain the empty type.
\end{definition}

\begin{theorem}[Substitution]
  \label{thm:aba-substitution}
  \lean{PLTS.ABA.ABDY.substSim}\lean{PLTS.ABA.ABDY.substitution}\leanok
  \uses{def:aba-hybrid, def:aba-composed, thm:gbca-family-refines, thm:precong-parallel, thm:precong-abstract, thm:precong-relabel, thm:forward-sound, def:achievable}
  Replacing each round's graded-agreement instance by its specification is a forward simulation of $\abdy{\mathrm{composed}}(P)$ by $\mathrm{hybrid}(P)$, \leandecl{PLTS.ABA.ABDY.substSim}, and hence $\operatorname{ATD}(\abdy{\mathrm{composed}}(P)) \subseteq \operatorname{ATD}(\mathrm{hybrid}(P))$: \leandecl{PLTS.ABA.ABDY.substitution}.
\end{theorem}

\begin{proof}\leanok
  \uses{def:aba-hybrid, def:aba-composed, thm:gbca-family-refines, thm:precong-parallel, thm:precong-abstract, thm:precong-relabel, thm:forward-sound}
  See the Lean proof of \leandecl{PLTS.ABA.ABDY.substitution}.
\end{proof}


\subsection{From the hybrid to the ABA specification}

The refinement of $\mathrm{hybrid}(P)$ into the specification is one hand-built simulation, run in the hybrid's own coordinates.
The round specifications, the $n$ round loops, the ABA network and the common coin are each a component of the system, and each is a component of the state the simulation relation is defined on.
The simulation's abstract state is lazy and two-phase (D16), deferring the decision to a single weak step under the first visible return, which is what a forward simulation permits and a strong one does not.
See \leanmodule{Leslie2Protocols/ABA/Composition/ABAState} and \leanmodule{Leslie2Protocols/ABA/HybridRefinesSpecification/Simulation}.

\begin{theorem}[The hybrid refines the ABA specification]
  \label{thm:aba-hybrid-spec}
  \lean{PLTS.ABA.HybridState}\lean{PLTS.ABA.Invariant}\lean{PLTS.ABA.AbstractState}\lean{PLTS.ABA.hybridSpecificationStateRelation}\lean{PLTS.ABA.hybridSpecificationRelation}\lean{PLTS.ABA.hybridRefinesSpecification}\lean{PLTS.ABA.hybrid_spec}\leanok
  \uses{def:aba-hybrid, def:aba-spec, def:aba-state, def:aba-gbca-spec, def:aba-wcc-spec, def:aba-labels, def:process-algebra, def:idle-family, def:forwardsim, def:weakstep, def:coupling, def:achievable, thm:forward-sound}
  The hybrid refines the ABA specification, $\operatorname{ATD}(\mathrm{hybrid}(P)) \subseteq \operatorname{ATD}(\mathrm{spec}(P))$, at every round message type: \leandecl{PLTS.ABA.hybrid_spec}.
\end{theorem}

\begin{proof}\leanok
  \uses{def:aba-hybrid, def:aba-spec, def:aba-state, def:aba-gbca-spec, def:aba-wcc-spec, def:aba-labels, def:process-algebra, def:idle-family, def:forwardsim, def:weakstep, def:coupling, def:achievable, thm:forward-sound}
  One hand-built probabilistic forward simulation, \leandecl{PLTS.ABA.hybridRefinesSpecification}, along the relation \leandecl{PLTS.ABA.hybridSpecificationRelation} on the state of $\mathrm{hybrid}(P)$ itself, then soundness (Theorem~\ref{thm:forward-sound}).
\end{proof}


\blueprintfig{fig-coresim}{The simulation of Theorem~\ref{thm:aba-hybrid-spec}, \leandecl{PLTS.ABA.hybridRefinesSpecification}: the two-phase abstract state, its stutter matching every hidden transition --- the coin's resolution, the one probabilistic transition among them, included --- under a constant coupling, and the single \leandeclnamed{decide}{PLTS.ABA.SpecStep.decide} step that leads the first $\mathrm{retABA}$, the two together forming the weak step that matches it.}

\subsection{The main theorems}

Three steps separate the protocol from the ABA specification:
\[
  \abdy{\mathrm{protocol}}(P)
  \;\overset{\text{(1)}}{\sqsubseteq}\; \abdy{\mathrm{composed}}(P)
  \;\overset{\text{(2)}}{\sqsubseteq}\; \mathrm{hybrid}(P)
  \;\overset{\text{(3)}}{\sqsubseteq}\; \mathrm{spec}(P),
\]
where (1) is the composition inclusion of Theorem~\ref{thm:aba-protocol-composed}, along the Dirac lift of $\mathrm{ProtocolRelation}(P)$; (2) the substitution of Theorem~\ref{thm:aba-substitution}; and (3) the refinement of Theorem~\ref{thm:aba-hybrid-spec}, whose proof is the hand-built core simulation, read in the coordinates of $\mathrm{hybrid}(P)$.
Safety then transfers by Theorem~\ref{thm:trace-support}, with no hypothesis on the trace beyond positive probability, and both properties are stated on the returns of never-corrupted processes.
The chain is assembled in \leanmodule{Leslie2Protocols/ABA/Results}.

\begin{theorem}[Safety of the protocol]
  \label{thm:aba-protocol}
  \lean{PLTS.ABA.hybrid_spec}\lean{PLTS.ABA.ABDY.protocol_safe}\lean{PLTS.ABA.ABDY.protocol_traces}\lean{PLTS.ABA.ABDY.composed_safe}\leanok
  \uses{def:aba-programs, def:aba-protocol, thm:aba-protocol-composed, thm:aba-substitution, thm:aba-hybrid-spec, thm:aba-spec-safe, thm:trace-support, thm:forward-sound, def:achievable}
  The protocol refines the ABA specification, and every trace of positive probability under an achievable trace distribution of it satisfies Validity and Agreement: \leandecl{PLTS.ABA.ABDY.refines}, \leandecl{PLTS.ABA.ABDY.protocol_safe}.
\end{theorem}

\begin{proof}\leanok
  \uses{thm:aba-protocol-composed, thm:aba-substitution, thm:aba-hybrid-spec, thm:aba-spec-safe, thm:trace-support, thm:forward-sound, def:achievable}
  See the Lean proofs of \leandecl{PLTS.ABA.ABDY.refines} and \leandecl{PLTS.ABA.ABDY.protocol_safe}.
\end{proof}

\begin{theorem}[Refinement and correctness of ABA]
  \label{thm:aba-main}
  \lean{PLTS.ABA.ABDY.refines}\lean{PLTS.ABA.ABDY.main}\leanok
  \uses{thm:aba-protocol, thm:aba-protocol-composed, thm:aba-substitution, thm:aba-hybrid-spec, thm:aba-spec-safe}
  Every trace distribution achievable by the protocol is achievable by the ABA specification, and every positive-probability trace of the protocol satisfies Validity and Agreement: \leandecl{PLTS.ABA.ABDY.refines}, \leandecl{PLTS.ABA.ABDY.main}.
\end{theorem}

\begin{proof}\leanok
  \uses{thm:aba-protocol, thm:aba-substitution, thm:aba-hybrid-spec, thm:aba-spec-safe, thm:trace-support, thm:forward-sound, thm:forward-trans}
  See the Lean proofs of \leandecl{PLTS.ABA.ABDY.refines} and \leandecl{PLTS.ABA.ABDY.main}.
\end{proof}


\subsection{Non-vacuity}

A refinement statement is compatible with a system that does nothing.
The chain is therefore accompanied by machine-checked witnesses that the hybrid of Definition~\ref{def:aba-hybrid}, the system Theorem~\ref{thm:aba-hybrid-spec} is proved of, executes and reaches a return at $n = 4$ and $f = 1$, on a connected execution of positive probability that exercises the one probabilistic transition on the way, and that the $\mathrm{fail}$ label fires in it rather than being merely declared.
Membership of a decision trace in the achievable trace distributions of the protocol is not established; the witnesses are about that system alone.
Both chains pass through that system, so the witnesses serve both.
See \leanmodule{Leslie2Protocols/ABA/HybridRefinesSpecification/NonVacuity}.

\begin{definition}[Non-vacuity]
  \label{def:aba-nonvacuity}
  \lean{PLTS.ABA.P4}\leanok
  \uses{thm:aba-hybrid-spec}
  Machine-checked witnesses that $\mathrm{hybrid}(P)$, the system Theorem~\ref{thm:aba-hybrid-spec} is proved of, executes and reaches a return on a connected run of positive probability, at the parameters $n = 4$ and $f = 1$ of \leandecl{PLTS.ABA.P4}: \leanmodule{Leslie2Protocols/ABA/Core/NonVacuity}.
  A synchronised $\mathrm{fail}\,0$ broadcast is witnessed separately, its network half discharging the two guards $0 \notin F$ and $|F| < f$ (D1) and its round-loop half writing the replacement flag (D23).
\end{definition}


\section{The gather-based implementation}
\label{sec:aba-gather}

Transition System~2 has a second verified implementation.
Attiya, Flam and Welch~\cite{AFW25} build graded agreement from two calls to a \emph{gather} sub-protocol and local counting, and the source blueprint carries the matching material: a gather specification (Transition System~4), a gather implementation over reliable broadcast (its Algorithm~4, the binding form of Algorithm~5 of~\cite{AFW25}), and Bracha's reliable broadcast (Transition System~6 and its Algorithm~6, from~\cite{Bra87}).
The Lean files are \leanmodule{Leslie2Protocols/ABA/Vocabulary/ProcessAndNetworkState} through \leanmodule{Leslie2Protocols/ABA/AFW/Substitution}, mapped file by file in the ABA library's own guide.

\subsection{Reliable broadcast}

Reliable broadcast sits at the bottom of both chains, and its specification is the first point dynamic corruption shapes: a leader corrupted after its call may still direct the delivered value, so the committed value is a separate field written by a guarded internal transition, not the field the call writes (D27).
The implementation is Bracha's three message levels over the shared two-part state of Definition~\ref{def:aba-networkstate}, and its refinement witnesses the committed value by a quorum of received $\mathrm{ECHO}$ messages, firing the abstract commit on demand under the first return.
See \leanmodule{Leslie2Protocols/ABA/Vocabulary/ProcessAndNetworkState}, \leanmodule{Leslie2Protocols/ABA/ReliableBroadcast/Specification}, \leanmodule{Leslie2Protocols/ABA/ReliableBroadcast/Bracha/MessagesAndVariables}, \leanmodule{Leslie2Protocols/ABA/ReliableBroadcast/Bracha/Components}, \leanmodule{Leslie2Protocols/ABA/ReliableBroadcast/Bracha/Composition}, \leanmodule{Leslie2Protocols/ABA/ReliableBroadcast/Bracha/SpecificationOverInstanceAlphabet}, \leanmodule{Leslie2Protocols/ABA/ReliableBroadcast/Bracha/CompositionStepCases}, \leanmodule{Leslie2Protocols/ABA/ReliableBroadcast/Bracha/Algorithm}, \leanmodule{Leslie2Protocols/ABA/ReliableBroadcast/Bracha/EchoWitness}, \leanmodule{Leslie2Protocols/ABA/ReliableBroadcast/Bracha/Invariant}, \leanmodule{Leslie2Protocols/ABA/ReliableBroadcast/Bracha/SpecificationRelation} and \leanmodule{Leslie2Protocols/ABA/ReliableBroadcast/Bracha/RefinesSpecification}.

\begin{definition}[The two-part instance vocabulary]
  \label{def:aba-networkstate}
  \lean{PLTS.ABA.NetworkState}\lean{PLTS.ABA.LocalState}\lean{PLTS.ABA.InstanceState}\lean{PLTS.ABA.InstanceState.corrupt}\lean{PLTS.ABA.InstanceState.multicast}\lean{PLTS.ABA.InstanceState.receiveMessage}\lean{PLTS.ABA.InstanceState.receivedCount}\leanok
  \uses{def:aba-parameters}
  The state shape the broadcast and gather encodings share: a network state holding the per-sender sent sets and the corrupted set (D5, D1) beside $n$ process-local states, each holding one process's variables and its per-sender received sets.
  Multicast, delivery, corruption and the counts of received messages are operations on the pair, and the quorum-intersection lemmas are stated once over it: \leandecl{PLTS.ABA.InstanceState}.
\end{definition}

\begin{definition}[The BRB specification instance]
  \label{def:aba-brb-spec}
  \lean{PLTS.ABA.BRB.Label}\lean{PLTS.ABA.BRB.SpecState}\lean{PLTS.ABA.BRB.SpecState.corrupt}\lean{PLTS.ABA.BRB.Step}\lean{PLTS.ABA.BRB.specInst}\leanok
  \uses{def:aba-params}
  Transition System~6, read safety-only over an instance-local alphabet, with the leader's recorded input and the committed value held in separate fields (D27): an internal commit writes the value once --- any value under a corrupted leader, the input otherwise --- and every return hands out the committed value.
  Totality, a liveness property, is not encoded: \leandecl{PLTS.ABA.BRB.specInst}.
\end{definition}


\blueprintfig{fig-brb-lifecycle}{The BRB specification instance \leandecl{PLTS.ABA.BRB.specInst} (Transition System 6): the reachable states quotiented to the pair of the leader's recorded input and the committed value.
The call records the leader's input, and the internal commit writes the delivered value once, under either of two guards (D27): the leader's own input, or any value once the leader is corrupted.
A corrupted leader therefore commits without a call at all, which is the edge leaving the initial state.
The guard of \leandeclnamed{commit}{PLTS.ABA.BRB.Step.commit} requires $val = \bot$ and no transition writes $val$ a second time, so the committed value is written once.
Every return hands out that value, so the returns are a self-loop on the committed state.
Totality is a liveness property and is not encoded; $\mathrm{fail}$ is omitted.}
\begin{definition}[The BRB implementation instance]
  \label{def:aba-bracha-implementation}
  \lean{PLTS.ABA.BRB.Message}\lean{PLTS.ABA.BRB.ProcessRecord}\lean{PLTS.ABA.BRB.BrachaState}\lean{PLTS.ABA.BRB.BroadcastEvent}\lean{PLTS.ABA.BRB.specificationLabelMap}\lean{PLTS.ABA.BRB.BrachaAlgorithm}\lean{PLTS.ABA.BRB.brachaInstance}\lean{PLTS.ABA.BRB.brachaInstance_step_iff_row}\leanok
  \uses{def:aba-brb-spec, def:aba-networkstate, def:synchronised-product, def:process-algebra, def:relabel}
  Bracha's algorithm~\cite{Bra87} --- Algorithm~1 of~\cite{AFW25} (D34) --- run as the $n$ per-process programs beside the instance's network, with the multicast and the delivery hidden: the leader multicasts $\mathrm{INIT}$, a process echoes once on the leader's $\mathrm{INIT}$ or on an $\mathrm{ECHO}$ quorum of more than $(n+f)/2$ senders or on $f{+}1$ $\mathrm{VOTE}$ receipts, votes once on that $\mathrm{ECHO}$ quorum or on $f{+}1$ $\mathrm{VOTE}$ receipts, and returns a value at $2f{+}1$ $\mathrm{VOTE}$ receipts of it --- \leandecl{PLTS.ABA.BRB.brachaInstance}.
  The specification answers its call label on two rows, the call and the input-enabledness loop; the composition takes the loop on a label of its own and reads the specification along \leandecl{PLTS.ABA.BRB.specificationLabelMap}, which sends the loop to the call.
  \leandecl{PLTS.ABA.BRB.brachaInstance_step_iff_row} states the transitions of the composition at one specification label as a list of cases, one per transition of \leandecl{PLTS.ABA.BRB.BrachaAlgorithm}.
\end{definition}

\begin{theorem}[BRB instance refinement]
  \label{thm:brb-refines}
  \lean{PLTS.ABA.BRB.EchoWitness}\lean{PLTS.ABA.BRB.Invariant}\lean{PLTS.ABA.BRB.SpecificationRelation}\lean{PLTS.ABA.BRB.specificationRelation_corrupt}\lean{PLTS.ABA.BRB.brachaRefinesSpecification}\leanok
  \uses{def:aba-bracha-implementation, def:aba-brb-spec}
  The Bracha instance is forward-simulated by Transition System~6, along a relation certifying the committed value by an $\mathrm{ECHO}$ receipt quorum of more than $(n+f)/2$ senders and firing the commit on demand under the first return; the relation survives the corruption broadcast: \leandecl{PLTS.ABA.BRB.brachaRefinesSpecification}, \leandecl{PLTS.ABA.BRB.specificationRelation_corrupt}.
\end{theorem}

\begin{proof}\leanok
  \uses{def:aba-bracha-implementation, def:aba-brb-spec, def:forwardsim}
  See the Lean proof of \leandecl{PLTS.ABA.BRB.brachaRefinesSpecification}.
\end{proof}


\subsection{Gather}

Gather has every process broadcast a payload and returns to each caller a set of (process, payload) entries, with a binding guarantee: the returns share a committed core fixed no later than the first of them.
The specification holds that core as write-once state of size at least $n - f$ and carries it on every return label, where the source carries a bound identifier set on the label alone.
The size bound is provable from the prefix of an execution, which is what the transition writing the core requires.
The core is one correct sender's $\mathrm{ECHO}$ payload, located on the sent sets by the counting of Definition~\ref{def:aba-gather-core}.
Remark~22 of~\cite{AFW25} determines the core of a non-binding gather only in hindsight; here it is an auxiliary variable in the sense of~\cite{AL91}, held as state and announced on the returns (D29).
The implementation is a composition: $n$ gather programs beside the gather network, which holds the core, in parallel with $2n$ reliable-broadcast instances read along pullbacks naming them, the instance's own events hidden.
A program keeps in its own variables what the broadcast instances have returned to it, written on the return events, and reads nothing else.
The implementation sends its $\mathrm{BIND}$ payloads by reliable broadcast, which is what fixes a payload against a sender corrupted later, and is verified in two substitutions: the Bracha instances replaced by congruence at every coordinate, then the whole instance replaced by Transition System~4 on the composition.
See \leanmodule{Leslie2Protocols/ABA/Gather/Specification}, \leanmodule{Leslie2Protocols/ABA/Gather/SpecificationSafety}, \leanmodule{Leslie2Protocols/ABA/Gather/MessagesAndCommonCore}, \leanmodule{Leslie2Protocols/ABA/Gather/Components}, \leanmodule{Leslie2Protocols/ABA/Gather/Composition}, \leanmodule{Leslie2Protocols/ABA/Gather/SpecificationOverInstanceAlphabet}, \leanmodule{Leslie2Protocols/ABA/Gather/CompositionStepCases}, \leanmodule{Leslie2Protocols/ABA/Gather/AlgorithmOverBracha}, \leanmodule{Leslie2Protocols/ABA/Gather/AlgorithmOverBroadcastSpecification}, \leanmodule{Leslie2Protocols/ABA/Gather/Invariant}, \leanmodule{Leslie2Protocols/ABA/Gather/CommonCoreCounting}, \leanmodule{Leslie2Protocols/ABA/Gather/SpecificationRelation}, \leanmodule{Leslie2Protocols/ABA/Gather/RefinesSpecification}, \leanmodule{Leslie2Protocols/ABA/Gather/BroadcastSubstitution} and \leanmodule{Leslie2Protocols/ABA/Gather/CommonCore}.

\begin{definition}[The gather specification instance]
  \label{def:aba-gather-spec}
  \lean{PLTS.ABA.Gather.AcceptedPairs}\lean{PLTS.ABA.Gather.AcceptedPairs.subMap}\lean{PLTS.ABA.Gather.Label}\lean{PLTS.ABA.Gather.SpecState}\lean{PLTS.ABA.Gather.SpecState.corrupt}\lean{PLTS.ABA.Gather.Step}\lean{PLTS.ABA.Gather.specInst}\leanok
  \uses{def:aba-parameters}
  Transition System~4 with the call/commit split and one bound core.
  Calls are recorded per process; each entry commits once, to any payload at a corrupted process and to the recorded call otherwise (D26); one internal transition writes, once, a set of committed entries of size at least $n-f$; and a return delivers a sub-map of the committed entries containing that set, which its label carries: \leandecl{PLTS.ABA.Gather.specInst}.
  The core is a ghost output, held as state and announced on every return, so a system built over the instance binds by reading its labels (D29).
\end{definition}


\blueprintfig{fig-gather-lifecycle}{The gather specification instance \leandecl{PLTS.ABA.Gather.specInst} (Transition System 4).
Above, the reachable states and transitions of one process's entry: the call records $k$'s payload, and \leandeclnamed{commit}{PLTS.ABA.Gather.Step.commit} writes $val(k)$ once, under $k$'s own call or, once $k$ is corrupted, under any value (D26).
The guard reads $F$ at the commit rather than at the call, so a process corrupted after its call may still direct its committed entry.
Below, the reachable states and transitions of the instance: \leandeclnamed{bindCore}{PLTS.ABA.Gather.Step.bindCore} writes the core $C$ once, requiring that its entries are already committed and that it has at least $n-f$ of them.
No return is enabled before that write, and a return delivers a sub-map of the committed entries containing $C$, which its label carries beside the returner and the map (D29).
$\mathrm{fail}$ is omitted.}
\begin{definition}[The gather instance as a composition]
  \label{def:aba-gather-composition}
  \lean{PLTS.ABA.Gather.ProcessVariables}\lean{PLTS.ABA.Gather.NetworkState}\lean{PLTS.ABA.Gather.GatherEvent}\lean{PLTS.ABA.Gather.inputBroadcastLabelMap}\lean{PLTS.ABA.Gather.bindBroadcastLabelMap}\lean{PLTS.ABA.Gather.specificationLabelMap}\lean{PLTS.ABA.Gather.StateOverBroadcasts}\lean{PLTS.ABA.Gather.instanceOverBroadcasts}\lean{PLTS.ABA.Gather.specificationOverInstanceAlphabet}\leanok
  \uses{def:aba-gather-spec, def:aba-bracha-implementation, def:aba-networkstate, def:synchronised-product, def:process-algebra, def:idle-family, def:relabel}
  One gather instance over a payload type $X$: the $n$ gather programs beside the gather network, in parallel with $2n$ reliable-broadcast instances, the instance's six events hidden and the result read back over its interface alphabet --- \leandecl{PLTS.ABA.Gather.instanceOverBroadcasts}, which takes the broadcast instances as an argument.
  A program holds its own variables, the messages delivered to it and two fields holding what the broadcasts returned, the value each input-broadcast instance and the payload each bind-broadcast instance has returned here, written on the return events and read by the guards that consume them; the gather network holds the per-sender sent sets, the corrupted set and the instance's core (D29).
  Each broadcast instance is read along a pullback that names it, \leandecl{PLTS.ABA.Gather.inputBroadcastLabelMap} or \leandecl{PLTS.ABA.Gather.bindBroadcastLabelMap}, so a label carrying another instance's index leaves that instance unchanged, and the specification is read along \leandecl{PLTS.ABA.Gather.specificationLabelMap}, which sends the call's input-enabledness loop to the call.
\end{definition}

\begin{definition}[The gather implementation over broadcast specifications]
  \label{def:aba-gather-over-broadcast-specification}
  \lean{PLTS.ABA.Gather.Message}\lean{PLTS.ABA.Gather.BaseProcessVariables}\lean{PLTS.ABA.Gather.StateOverBroadcastSpecification}\lean{PLTS.ABA.Gather.AlgorithmOverBroadcastSpecification}\lean{PLTS.ABA.Gather.instanceOverBroadcastSpecification}\lean{PLTS.ABA.Gather.instanceOverBroadcastSpecification_step_iff_algorithm}\leanok
  \uses{def:aba-gather-composition, def:aba-gather-spec, def:aba-brb-spec, def:aba-networkstate}
  The gather instance of Definition~\ref{def:aba-gather-composition} whose broadcast tier is $2n$ broadcast specifications: \leandecl{PLTS.ABA.Gather.instanceOverBroadcastSpecification}.
  Its transitions are Algorithm~4 of~\cite{LeslieBP}, the binding form of Algorithm~5 of~\cite{AFW25}, stated over the composition's state: each process contributes its input and its $\mathrm{BIND}$ payload through its own broadcast instance (D32), an entry set is approved at a process when that process's input broadcasts have returned every entry of it, a process multicasts the $n{-}f$ or more entries its own input broadcasts returned as its $\mathrm{ECHO}$, and the $\mathrm{ECHO}$/$\mathrm{VOTE}$ exchange over approved sets runs on the gather network.
  Each send after the $\mathrm{ECHO}$ requires the sender's own send at the level below, and the bind call requires that no earlier bind call stands in the write-once field $\mathit{sentBind}$ of the record (D35).
  A return requires the returner's own bind call, $n{-}f$ bind stores below the output and the output below the input stores, and the gather network writes the core at the first of them (D29): \leandecl{PLTS.ABA.Gather.AlgorithmOverBroadcastSpecification}, matched against the transitions of the composition by \leandecl{PLTS.ABA.Gather.instanceOverBroadcastSpecification_step_iff_algorithm}.
\end{definition}

\begin{definition}[The core of a gather message state]
  \label{def:aba-gather-core}
  \lean{PLTS.ABA.Gather.honest}\lean{PLTS.ABA.Gather.dominatedBy}\lean{PLTS.ABA.Gather.dominators}\lean{PLTS.ABA.Gather.echoOf}\lean{PLTS.ABA.Gather.coreOf}\lean{PLTS.ABA.Gather.IdealInv}\lean{PLTS.ABA.Gather.single_core}\lean{PLTS.ABA.Gather.single_core_approved}\lean{PLTS.ABA.Gather.bindAbove}\lean{PLTS.ABA.Gather.coreOf_freeze}\leanok
  \uses{def:aba-gather-ideal, def:aba-params}
  A payload set read off the sent sets and the corrupted set of a gather message state: the $\mathrm{ECHO}$ payload of a sender whose payload lies below every $\mathrm{VOTE}$ of many processes outside $F$ (\leandecl{PLTS.ABA.Gather.coreOf}).
  Under the invariant \leandecl{PLTS.ABA.Gather.IdealInv}, as soon as one process outside $F$ holds a committed $\mathrm{BIND}$ payload, that set has at least $n-f$ entries, its entries are committed input-broadcast entries, and it lies below the committed $\mathrm{BIND}$ payload of every process outside $F$: \leandecl{PLTS.ABA.Gather.single_core} and \leandecl{PLTS.ABA.Gather.single_core_approved}.
  The count of coordinates holding a committed $\mathrm{BIND}$ payload above it is $F$-blind and monotone, and an $n{-}f$ return quorum supplies $f{+}1$ of them: \leandecl{PLTS.ABA.Gather.bindAbove} and \leandecl{PLTS.ABA.Gather.coreOf_freeze}.
\end{definition}

\begin{theorem}[Gather core refinement]
  \label{thm:gather-core}
  \lean{PLTS.ABA.Gather.Invariant}\lean{PLTS.ABA.Gather.SpecificationRelation}\lean{PLTS.ABA.Gather.refinesSpecification}\leanok
  \uses{def:aba-gather-over-broadcast-specification, def:aba-gather-spec, def:aba-gather-core}
  The gather implementation over broadcast specifications is forward-simulated by Transition System~4: \leandecl{PLTS.ABA.Gather.refinesSpecification}.
  The core the abstract system holds is the core of the gather network's state, and a return is matched by one weak run that commits the abstract entries, writes the core where it is unwritten, and returns.
  The refinement is a case analysis over the transitions of Definition~\ref{def:aba-gather-over-broadcast-specification}, through its characterisation by the algorithm; the specification's call record follows the input instances' records, whose writes it shares.
\end{theorem}

\begin{proof}\leanok
  \uses{def:aba-gather-over-broadcast-specification, def:aba-gather-spec, def:aba-gather-core, def:forwardsim}
  See the Lean proof of \leandecl{PLTS.ABA.Gather.refinesSpecification}.
\end{proof}

\begin{definition}[The gather implementation over Bracha instances]
  \label{def:aba-gather-over-bracha}
  \lean{PLTS.ABA.Gather.StateOverBracha}\lean{PLTS.ABA.Gather.AlgorithmOverBracha}\lean{PLTS.ABA.Gather.instanceOverBracha}\lean{PLTS.ABA.Gather.instanceOverBracha_step_iff_algorithm}\leanok
  \uses{def:aba-gather-composition, def:aba-gather-over-broadcast-specification, def:aba-bracha-implementation}
  The gather instance of Definition~\ref{def:aba-gather-composition} whose $2n$ broadcast coordinates are Bracha instances of Definition~\ref{def:aba-bracha-implementation}: \leandecl{PLTS.ABA.Gather.instanceOverBracha}.
  Its transitions are those of Definition~\ref{def:aba-gather-over-broadcast-specification} with each broadcast transition carrying a Bracha transition in its place, and the internal transitions of the $2n$ instances --- their deliveries, $\mathrm{INIT}$ multicasts, echoes, votes and injections --- are the silent transitions of the composition: \leandecl{PLTS.ABA.Gather.AlgorithmOverBracha}, matched against the transitions of the composition by \leandecl{PLTS.ABA.Gather.instanceOverBracha_step_iff_algorithm}.
  The gather programs and the gather network are the same in both compositions, so the fields holding what the broadcasts returned, the guards that read them and the core are written here as they are over the broadcast specifications (D29).
\end{definition}

\begin{theorem}[Gather broadcast substitution]
  \label{thm:gather-broadcast-substitution}
  \lean{PLTS.ABA.Gather.BroadcastSubstitutionRelation}\lean{PLTS.ABA.Gather.broadcastSubstitution}\leanok
  \uses{def:aba-gather-over-bracha, def:aba-gather-over-broadcast-specification, thm:brb-refines, thm:forwardsim-congruence}
  The gather implementation over Bracha instances is forward-simulated by the gather implementation over broadcast specifications, by congruence: the broadcast refinement of Theorem~\ref{thm:brb-refines}, read along the pullback naming each of the $2n$ coordinates, is carried through the two products, the parallel composition with the gather programs and the gather network held, the hiding and the restriction by Theorem~\ref{thm:forwardsim-congruence} --- \leandecl{PLTS.ABA.Gather.broadcastSubstitution}, along \leandecl{PLTS.ABA.Gather.BroadcastSubstitutionRelation}, which holds those two components equal beside the broadcast relation at every coordinate.
\end{theorem}

\begin{proof}\leanok
  \uses{def:aba-gather-over-bracha, def:aba-gather-over-broadcast-specification, thm:brb-refines, def:forwardsim}
  See the Lean proof of \leandecl{PLTS.ABA.Gather.broadcastSubstitution}.
\end{proof}

\begin{theorem}[The common core on a trace]
  \label{thm:aba-gather-commoncore}
  \lean{PLTS.ABA.Gather.CoreTrace}\lean{PLTS.ABA.Gather.specInst_core}\lean{PLTS.ABA.Gather.instanceOverBroadcastSpecification_core}\lean{PLTS.ABA.Gather.instanceOverBracha_core}\leanok
  \uses{def:aba-gather-spec, def:aba-gather-over-broadcast-specification, def:aba-gather-over-bracha, thm:gather-core, thm:gather-broadcast-substitution, thm:trace-support}
  Every trace of positive probability of the gather specification instance satisfies $\mathrm{CoreTrace}$: all its return labels carry one and the same payload set, that set has at least $n-f$ entries, and each returned map has every entry of it (\leandecl{PLTS.ABA.Gather.CoreTrace}, \leandecl{PLTS.ABA.Gather.specInst_core}).
  Both implementations inherit it along their refinements: \leandecl{PLTS.ABA.Gather.instanceOverBroadcastSpecification_core} and \leandecl{PLTS.ABA.Gather.instanceOverBracha_core}.
\end{theorem}

\begin{proof}\leanok
  \uses{def:aba-gather-spec, thm:gather-core, thm:gather-broadcast-substitution, thm:trace-support, thm:forward-sound}
  The specification clauses come from the transitions, and the two implementations inherit them along the soundness of their refinements; see \leandecl{PLTS.ABA.Gather.specInst_core}.
\end{proof}


\subsection{The two-gather round}

The round calls gather on the input bits, takes as candidate the bit carried by all but $f$ of the returned entries, calls gather again on the candidates, and reads the grade off the counts of the second return (D24).
Both gathers are the binding form (D33).
The round is a composition too: $n$ programs beside the round's own network, which carries no messages and holds the round's bound bit, in parallel with the two gathers read along pullbacks naming them.
The first gather's return to a process and that process's call of the second gather are two events, and so are the second gather's return and the round's graded return; the program's variables carry the round across each of them.
The refinement into Transition System~2 is pure counting on the two cores, each written once: a bit the first core carries on fewer than $|S| - f$ entries is excluded and stays excluded, a core carrying one candidate on at least $|S| - f$ entries leaves the other on fewer, and the grade-2 / grade-0 guard is the second core carrying a bit on at least $|S| - f$ entries against its carrying each bit on at most $f$.
The round's bound bit is the bit the first core carries on at least $|S| - f$ entries, written at the first gather's first return and announced by every graded return, and it is the bit whose complement each return excludes.
See \leanmodule{Leslie2Protocols/ABA/GBCA/AFW/Counting}, \leanmodule{Leslie2Protocols/ABA/GBCA/AFW/Components}, \leanmodule{Leslie2Protocols/ABA/GBCA/AFW/Composition}, \leanmodule{Leslie2Protocols/ABA/GBCA/AFW/CompositionStepCases}, \leanmodule{Leslie2Protocols/ABA/GBCA/AFW/AlgorithmOverGatherSpecifications}, \leanmodule{Leslie2Protocols/ABA/GBCA/AFW/OutputWitness}, \leanmodule{Leslie2Protocols/ABA/GBCA/AFW/Invariant}, \leanmodule{Leslie2Protocols/ABA/GBCA/AFW/SpecificationRelation}, \leanmodule{Leslie2Protocols/ABA/GBCA/AFW/RefinesSpecification}, \leanmodule{Leslie2Protocols/ABA/GBCA/AFW/BroadcastSubstitution}, \leanmodule{Leslie2Protocols/ABA/GBCA/AFW/GatherSubstitution} and \leanmodule{Leslie2Protocols/ABA/GBCA/AFW/Binding}.

\begin{definition}[The round's network]
  \label{def:aba-gbca-network}
  \lean{PLTS.ABA.GBCA.ByAFW.NetworkStep}\lean{PLTS.ABA.GBCA.ByAFW.GBCANetwork}\leanok
  \uses{def:aba-gather-spec, def:aba-labels}
  The network component of a gather-based round: it exchanges no messages and its whole state is the round's bound bit --- \leandecl{PLTS.ABA.GBCA.ByAFW.GBCANetwork}.
  The first gather's return writes the bit, where the component holds none, as $\mathrm{boundOfCore}$ of the core that return carries; every graded return announces the bit on its label and moves nothing; every other label leaves it where it stands.
  No program reads it, so it is auxiliary state (D29).
\end{definition}

\begin{definition}[The graded-agreement round as a composition]
  \label{def:aba-round-composition}
  \lean{PLTS.ABA.GBCA.ByAFW.ProcessVariables}\lean{PLTS.ABA.GBCA.ByAFW.RoundEvent}\lean{PLTS.ABA.GBCA.ByAFW.ProgramLabel}\lean{PLTS.ABA.GBCA.ByAFW.programLabelMap}\lean{PLTS.ABA.GBCA.ByAFW.firstGatherLabelMap}\lean{PLTS.ABA.GBCA.ByAFW.secondGatherLabelMap}\lean{PLTS.ABA.GBCA.ByAFW.roundPrograms}\lean{PLTS.ABA.GBCA.ByAFW.RoundStateOverGathers}\lean{PLTS.ABA.GBCA.ByAFW.roundOverGathers}\leanok
  \uses{def:aba-gather-composition, def:aba-gbca-network, def:aba-labels, def:synchronised-product, def:process-algebra, def:idle-family, def:relabel}
  The round-$r$ graded-agreement round over two gather instances: the $n$ graded-agreement programs beside the network of Definition~\ref{def:aba-gbca-network}, in parallel with two gather instances of Definition~\ref{def:aba-gather-composition}, the round's three events hidden and the result read over the family alphabet $\mathrm{ExtendedLabel}(n, M)$ at the empty round message type --- \leandecl{PLTS.ABA.GBCA.ByAFW.roundOverGathers}, which takes the two gather systems as arguments.
  A program holds one process's record of the round --- its input, its candidate, whether it has called the second gather, its graded outcome and its return flag --- and moves at each of the round's four moments: the first gather's return and the second gather's call are two events, and so are the second gather's return and the round's graded return.
  Each gather is read along a pullback that names it, \leandecl{PLTS.ABA.GBCA.ByAFW.firstGatherLabelMap} or \leandecl{PLTS.ABA.GBCA.ByAFW.secondGatherLabelMap}, and each program along \leandecl{PLTS.ABA.GBCA.ByAFW.programLabelMap}, which reads a Byzantine call as a call, a Byzantine return as a return, and the family's two call loops as the loop.
\end{definition}

\begin{definition}[The two-gather round]
  \label{def:aba-round-over-gather-specifications}
  \lean{PLTS.ABA.GBCA.candidate}\lean{PLTS.ABA.GBCA.gradeOf}\lean{PLTS.ABA.GBCA.boundOfCore}\lean{PLTS.ABA.GBCA.ByAFW.RoundStateOverGatherSpecifications}\lean{PLTS.ABA.GBCA.ByAFW.AlgorithmOverGatherSpecifications}\lean{PLTS.ABA.GBCA.ByAFW.roundOverGatherSpecifications}\lean{PLTS.ABA.GBCA.ByAFW.roundOverGatherSpecifications_step_iff_algorithm}\leanok
  \uses{def:aba-round-composition, def:aba-gbca-network, def:aba-gather-spec, def:aba-gbca-spec, def:aba-labels}
  The graded agreement of~\cite{AFW25} (its Algorithm~4 at $R = 2$, the two-gather branch, its approximate-agreement subroutine replaced by a local count, D24): the round of Definition~\ref{def:aba-round-composition} over two gather specification instances, each the binding form (D33), and the algorithm listing its transitions at the specification's labels --- \leandecl{PLTS.ABA.GBCA.ByAFW.roundOverGatherSpecifications}, \leandecl{PLTS.ABA.GBCA.ByAFW.AlgorithmOverGatherSpecifications}, matched against the composition by \leandecl{PLTS.ABA.GBCA.ByAFW.roundOverGatherSpecifications_step_iff_algorithm}.
  The round's call is the first gather's call; when the first gather returns $g$ the caller records the candidate --- the bit on all but $f$ of $g$'s entries, $\bot$ if neither bit is --- and calls the second gather with it; when the second gather returns $h$ the caller records the grade, $(v, 2)$ at $|{\rm dom}\, h| - f$ entries of $v$, $(v, 1)$ at $f{+}1$ and $(\bot, 0)$ below both, and the round returns it.
  The network of Definition~\ref{def:aba-gbca-network} writes the round's bound bit at the first gather's first return, as \leandecl{PLTS.ABA.GBCA.boundOfCore} of the core that return carries, and every graded return announces it (D29).
\end{definition}

\begin{theorem}[Two-gather round refinement]
  \label{thm:gbca-round-refines-specification}
  \lean{PLTS.ABA.GBCA.count_aboveThreshold_of_subMap}\lean{PLTS.ABA.GBCA.boundOfCore_of_aboveThreshold}\lean{PLTS.ABA.GBCA.count_boundOfCore_belowThreshold}\lean{PLTS.ABA.GBCA.ByAFW.Invariant}\lean{PLTS.ABA.GBCA.ByAFW.ExclusionWitness}\lean{PLTS.ABA.GBCA.ByAFW.SpecificationRelation}\lean{PLTS.ABA.GBCA.ByAFW.refinesSpecification}\leanok
  \uses{def:aba-round-over-gather-specifications, def:aba-gbca-spec}
  The two-gather round is forward-simulated by the GBCA specification instance (Transition System~2): \leandecl{PLTS.ABA.GBCA.ByAFW.refinesSpecification}.
  The exclusion set and the grade guard are certified on the two cores, each written once, --- a bit is excluded when the first core carries it on fewer than $|S| - f$ entries --- and excludes fire on demand under the returns that need them.
  The bit a return announces is the round's bound bit, so the specification's guard on the announcement and its guard pair on the value are the same pair.
  The candidate's certificate is recorded at the first gather's return and consumed at the second gather's call, and the grade's at the second gather's return and consumed at the graded return, the program's record carrying each across its two events.
\end{theorem}

\begin{proof}\leanok
  \uses{def:aba-round-over-gather-specifications, def:aba-gbca-spec, def:forwardsim}
  See the Lean proof of \leandecl{PLTS.ABA.GBCA.ByAFW.refinesSpecification}.
\end{proof}

\begin{definition}[The gather-based GBCA implementation]
  \label{def:aba-round-over-bracha}
  \lean{PLTS.ABA.GBCA.ByAFW.RoundStateOverBroadcastSpecification}\lean{PLTS.ABA.GBCA.ByAFW.roundOverBroadcastSpecification}\lean{PLTS.ABA.GBCA.ByAFW.RoundStateOverBracha}\lean{PLTS.ABA.GBCA.ByAFW.roundOverBracha}\leanok
  \uses{def:aba-round-composition, def:aba-round-over-gather-specifications, def:aba-gather-over-broadcast-specification, def:aba-gather-over-bracha}
  The round of Definition~\ref{def:aba-round-composition} at two concrete gather systems in its two gather positions: over gather instances of Definition~\ref{def:aba-gather-over-broadcast-specification}, and fully concrete over gather instances of Definition~\ref{def:aba-gather-over-bracha} --- the two gather instances and, beneath them, the $4n$ Bracha instances carrying the inputs and the $\mathrm{BIND}$ payloads.
  The concrete composition \leandecl{PLTS.ABA.GBCA.ByAFW.roundOverBracha} is the gather-based GBCA implementation, the protocol as it runs; \leandecl{PLTS.ABA.GBCA.ByAFW.roundOverBroadcastSpecification} is the composition between the two substitutions.
\end{definition}


\blueprintfig{fig-gather-compositions}{One round of the gather-based implementation, \leandecl{PLTS.ABA.GBCA.ByAFW.roundOverBracha}, drawn open, and the three substitutions that carry it to Transition System~2.
A round is $n$ graded-agreement programs beside the round's own network, which holds the bound bit and nothing else, in parallel with two gather instances of the binding form (D33), one over the input bits and one over the candidates; each gather is $n$ programs beside the gather network, which holds the core, in parallel with $2n$ Bracha instances.
That is $4n + 2$ networks in a single round --- $2n$ broadcasts for the inputs and $2n$ for the $\mathrm{BIND}$ payloads (D32, D33), beside the two gather networks --- which is the count the implementation collapses into one tagged sent-set family.
The three systems above it --- the three that carry one round, from the two gather specifications down to Bracha --- replace one component at a time: \leandecl{PLTS.ABA.GBCA.ByAFW.broadcastSubstitution} replaces each Bracha instance by the BRB specification, by congruence, \leandecl{PLTS.ABA.GBCA.ByAFW.gatherSubstitution} each gather instance by the gather specification, by congruence, and \leandecl{PLTS.ABA.GBCA.ByAFW.refinesSpecification} reads the grade off the two cores, each written once, and the exclusion set off the round's bound bit.
In each of them the round's own network holds that bound bit, written at the first gather's first return and untouched by either substitution.
Each is a plain forward simulation, so the three compose.}
\begin{theorem}[The two substitutions inside the round]
  \label{thm:gbca-round-substitutions}
  \lean{PLTS.ABA.GBCA.ByAFW.BroadcastSubstitutionRelation}\lean{PLTS.ABA.GBCA.ByAFW.broadcastSubstitution}\lean{PLTS.ABA.GBCA.ByAFW.GatherSubstitutionRelation}\lean{PLTS.ABA.GBCA.ByAFW.gatherSubstitution}\leanok
  \uses{def:aba-round-over-bracha, def:aba-round-over-gather-specifications, thm:gather-broadcast-substitution, thm:gather-core, thm:forwardsim-congruence}
  Each composition of the round is forward-simulated by the one above it, by congruence: the gather simulation of Theorem~\ref{thm:gather-broadcast-substitution}, respectively Theorem~\ref{thm:gather-core}, is read along the pullback naming each of the two gather positions and carried through the round's operators with the round's programs and network held, by Theorem~\ref{thm:forwardsim-congruence} --- \leandecl{PLTS.ABA.GBCA.ByAFW.broadcastSubstitution} along \leandecl{PLTS.ABA.GBCA.ByAFW.BroadcastSubstitutionRelation} and \leandecl{PLTS.ABA.GBCA.ByAFW.gatherSubstitution} along \leandecl{PLTS.ABA.GBCA.ByAFW.GatherSubstitutionRelation}, each holding the round's programs and network equal beside the gather relation at both positions.
\end{theorem}

\begin{proof}\leanok
  \uses{def:aba-round-over-bracha, def:aba-round-over-gather-specifications, thm:gather-broadcast-substitution, thm:gather-core, def:forwardsim}
  See the Lean proofs of \leandecl{PLTS.ABA.GBCA.ByAFW.broadcastSubstitution} and \leandecl{PLTS.ABA.GBCA.ByAFW.gatherSubstitution}.
\end{proof}


\blueprintfig{fig-gather-sim}{The three simulations of one round, with the concrete transitions on the right and the relation dashed.
Each stage idealises away one component and holds the components above it equal.
Under \leandecl{PLTS.ABA.GBCA.ByAFW.BroadcastSubstitutionRelation} every $\mathrm{ECHO}$, $\mathrm{VOTE}$ and delivery of a Bracha instance is matched by a stutter of the specification coordinate that replaces it.
Under \leandecl{PLTS.ABA.GBCA.ByAFW.GatherSubstitutionRelation} gather's own transitions stutter in the same way, and a broadcast return, which writes a program's variables, is matched by the specification's return.
Under \leandecl{PLTS.ABA.GBCA.ByAFW.SpecificationRelation} the return transitions carry the refinement: a return whose bit the specification has not yet excluded is matched by the weak run $\mathrm{bindUnset};\mathrm{ret}$, and a return whose exclusion has already happened is a single graded return.
The round's own network is held equal by the two substitutions, and \leandecl{PLTS.ABA.GBCA.ByAFW.SpecificationRelation} holds $excluded$ equal to the complement of the bit standing there.}
\begin{theorem}[Gather-based GBCA round refinement]
  \label{thm:gbca-gather-refines}
  \lean{PLTS.ABA.GBCA.roundOverBracha_refinesSpecification}\lean{PLTS.ABA.GBCA.roundOverBracha_specificationTraces}\leanok
  \uses{thm:gbca-tier-refines, thm:gbca-pair-refines, thm:forward-trans, lem:hyperstep-weakstep}
  The gather-based GBCA implementation refines the GBCA specification instance, round by round: the three simulations, taken probabilistically, compose to \leandecl{PLTS.ABA.GBCA.roundOverBracha_refinesSpecification}, and $\operatorname{ATD}$ inclusion follows: \leandecl{PLTS.ABA.GBCA.roundOverBracha_specificationTraces}.
\end{theorem}

\begin{proof}\leanok
  \uses{thm:gbca-tier-refines, thm:gbca-pair-refines, thm:forward-trans, lem:hyperstep-weakstep}
  See the Lean proof of \leandecl{PLTS.ABA.GBCA.roundOverBracha_refinesSpecification}.
\end{proof}


\begin{theorem}[Binding of graded agreement, on a trace]
  \label{thm:aba-gbca-binding-trace}
  \lean{PLTS.ABA.GBCA.BoundTrace}\lean{PLTS.ABA.GBCA.BindingTrace}\lean{PLTS.ABA.GBCA.specInst_binding}\lean{PLTS.ABA.GBCA.BindingTrace.value_agree}\lean{PLTS.ABA.GBCA.ByABDY.composition_binding}\lean{PLTS.ABA.GBCA.BindingTraceExtended}\lean{PLTS.ABA.GBCA.roundOverGatherSpecifications_binding}\lean{PLTS.ABA.GBCA.roundOverBroadcastSpecification_binding}\lean{PLTS.ABA.GBCA.roundOverBracha_binding}\leanok
  \uses{def:aba-gbca-spec, thm:gbca-spec-binding, thm:trace-support, thm:forward-sound, thm:aba-gbca-composition-by-abdy-refines-specification, thm:gbca-round-refines-specification, thm:gbca-round-substitutions, thm:gbca-gather-refines}
  Every trace of positive probability of the round-$r$ GBCA specification instance is bound to one bit: all its round-$r$ returns announce the same bit, and every one of them that hands out a value hands out that bit (\leandecl{PLTS.ABA.GBCA.BindingTrace}, \leandecl{PLTS.ABA.GBCA.specInst_binding}).
  Graded agreement follows on the trace alone: any two round-$r$ returns of it that hand out a bit hand out the same bit (\leandecl{PLTS.ABA.GBCA.BindingTrace.value_agree}).
  The property is stated on labels, so it transfers along the inclusions of Theorems~\ref{thm:aba-gbca-composition-by-abdy-refines-specification}, \ref{thm:gbca-round-refines-specification}, \ref{thm:gbca-round-substitutions} and~\ref{thm:gbca-gather-refines} to the direct implementation instance and --- read through the family pullback as \leandecl{PLTS.ABA.GBCA.BindingTraceExtended}, since the gather-based compositions are over the family alphabet --- to the round over the gather specifications, the round over the broadcast specifications and the gather-based implementation, in \leandecl{PLTS.ABA.GBCA.ByABDY.composition_binding}, \leandecl{PLTS.ABA.GBCA.roundOverGatherSpecifications_binding}, \leandecl{PLTS.ABA.GBCA.roundOverBroadcastSpecification_binding} and \leandecl{PLTS.ABA.GBCA.roundOverBracha_binding}.
\end{theorem}

\begin{proof}\leanok
  \uses{def:aba-gbca-spec, thm:gbca-spec-binding, thm:trace-support, thm:forward-sound, thm:aba-gbca-composition-by-abdy-refines-specification, thm:gbca-round-refines-specification, thm:gbca-round-substitutions, thm:gbca-gather-refines}
  See the Lean proof of \leandecl{PLTS.ABA.GBCA.specInst_binding}, and the four transfers along it in \leandecl{PLTS.ABA.GBCA.ByABDY.composition_specificationTraces}, \leandecl{PLTS.ABA.GBCA.roundOverGatherSpecifications_refines}, \leandecl{PLTS.ABA.GBCA.roundOverBroadcastSpecification_specificationTraces} and \leandecl{PLTS.ABA.GBCA.roundOverBracha_specificationTraces}.
\end{proof}


\section{The gather-based chain}

Beneath the composed system is the protocol that runs.
It is Definition~\ref{def:aba-implementation} a second time, at the gather-based implementation: the same round loop, the same DECIDED sets, the same coin's call and return, the same corruption discipline and the same pipeline, over a graded-agreement component that is the two gather instances and the $4n$ Bracha instances beneath them.
Three things separate that component from the composed system's.
A round of the composed system carries $4n+2$ networks; the adversary of the implementation carries one sent-set family, so a round message is tagged with the network it belongs to.
The composed system indexes the local states by instance and then by process, where a program must hold its own data and no one else's, so the round variables are transposed to hold one process's local state in every instance.
And a program of the implementation takes one transition per statement of the pseudocode: the graded-agreement call records the round's input, the call of a broadcast instance and the leader's multicast of $\langle\mathrm{INIT}, x\rangle$ that follows it are two, a gather's return and the call that follows it are two, and a broadcast instance's return is the transition that writes what it returned.
Assembling the composed state from the implementation's undoes all three, which is what the simulation relation of Theorem~\ref{thm:aba-afw-protocol-composed} does.
The adversary's ghost for a round is the third thing this instantiation supplies (D30): the two gathers' cores beside the round's bound bit, each written once.
The first gather's return writes the first core, off the projection of the tagged sent sets onto the first gather, and the bound bit as the bit that core carries on at least $|S| - f$ entries; the second gather's return writes the second core the same way, and a graded return announces the bit.
The composed round holds those three in its network components, so the relation reads them off the ghost and a transition's ghost write is the round's write of them.
See \leanmodule{Leslie2Protocols/ABA/AFW/System}, \leanmodule{Leslie2Protocols/ABA/AFW/RoundProjection}, \leanmodule{Leslie2Protocols/ABA/AFW/RoundProjectionStep}, \leanmodule{Leslie2Protocols/ABA/AFW/SimulationOfEachTransition} and \leanmodule{Leslie2Protocols/ABA/AFW/Simulation}.
\leanmodule{Leslie2Protocols/ABA/AFW/RoundProjectionStep} holds \leanmodule{Leslie2Protocols/ABA/AFW/RoundProjectionStep/ProjectionAfterOneWrite}, \leanmodule{Leslie2Protocols/ABA/AFW/RoundProjectionStep/GatherSend}, \leanmodule{Leslie2Protocols/ABA/AFW/RoundProjectionStep/BroadcastSend}, \leanmodule{Leslie2Protocols/ABA/AFW/RoundProjectionStep/BroadcastReturn}, \leanmodule{Leslie2Protocols/ABA/AFW/RoundProjectionStep/GatherAndBroadcastTransitions}, \leanmodule{Leslie2Protocols/ABA/AFW/RoundProjectionStep/GradedAgreementCall}, \leanmodule{Leslie2Protocols/ABA/AFW/RoundProjectionStep/FirstGatherReturn}, \leanmodule{Leslie2Protocols/ABA/AFW/RoundProjectionStep/SecondGatherCall}, \leanmodule{Leslie2Protocols/ABA/AFW/RoundProjectionStep/SecondGatherReturn}, \leanmodule{Leslie2Protocols/ABA/AFW/RoundProjectionStep/GradedReturn}, \leanmodule{Leslie2Protocols/ABA/AFW/RoundProjectionStep/Delivery}, \leanmodule{Leslie2Protocols/ABA/AFW/RoundProjectionStep/ByzantineInjection}, \leanmodule{Leslie2Protocols/ABA/AFW/RoundProjectionStep/OtherRoundsUnchanged} and \leanmodule{Leslie2Protocols/ABA/AFW/RoundProjectionStep/ProtocolRelationConjuncts}.

\begin{definition}[The gather-based protocol]
  \label{def:aba-afw-protocol}
  \lean{PLTS.ABA.AFW.Message}\lean{PLTS.ABA.AFW.RoundVariables}\lean{PLTS.ABA.AFW.RoundStep}\lean{PLTS.ABA.AFW.gbcaCallPayload}\lean{PLTS.ABA.AFW.Ghost}\lean{PLTS.ABA.AFW.firstGatherOf}\lean{PLTS.ABA.AFW.secondGatherOf}\lean{PLTS.ABA.AFW.ghostStep}\lean{PLTS.ABA.AFW.ghostOutput}\lean{PLTS.ABA.AFW.protocolExtended}\lean{PLTS.ABA.AFW.protocolHidden}\lean{PLTS.ABA.AFW.protocol}\leanok
  \uses{def:aba-implementation, def:aba-round-over-bracha, def:aba-gather-over-bracha, def:aba-gather-core}
  Definition~\ref{def:aba-implementation} at the gather-based implementation: one tagged message type collapsing a round's $4n+2$ network states into one sent-set family, process-major round variables holding each process's local state in every instance, and 31 transitions of the round --- the round's own calls and returns, the two gather instances, the $4n$ Bracha instances beneath them and the delivery: \leandecl{PLTS.ABA.AFW.protocol}.
  The adversary's ghost for a round is \leandecl{PLTS.ABA.AFW.Ghost}, the two gathers' cores and the round's bound bit, each written once: the first gather's return writes the first core and the bit, and the second gather's return writes the second (D29, D30).
\end{definition}


\blueprintfig{fig-afw-protocol}{The gather-based protocol \leandecl{PLTS.ABA.AFW.protocol}: the $n$ programs beside the one network and the common coin, under the two hidings of Definition~\ref{def:aba-implementation}.
The shape is that of Definition~\ref{def:aba-implementation}, and two things distinguish it from ABDY22's instantiation.
The adversary holds a single sent-set family per round over the tagged message type \leandecl{PLTS.ABA.AFW.Message}, whose tag names the network state a message belongs to and, for a broadcast message, the instance; the composed system of the same round carries $4n + 2$ separate network states.
And a program's variables are process-major: process $j$ holds its own local state in every instance, where the composed system indexes the local states by instance and then by process.
The adversary further holds the ghost of each round, \leandecl{PLTS.ABA.AFW.Ghost}: the two gathers' cores and the round's bound bit, each written once (D30).
Assembling the composed state from the implementation's undoes both rearrangements, which is what \leandecl{PLTS.ABA.AFW.roundProjection} does --- projecting the one sent-set family onto the network states, transposing the variables back, and reading the two cores and the bound bit off that ghost.}

At the protocol shape, the gather-based rounds occupy the same refinable position the direct rounds do: each round is over the extended alphabet itself, is gathered into the round-indexed family, and is put through the composed system's pipeline.
The three substitutions inside the round become three stages between protocol-shaped systems, the third landing on $\mathrm{hybrid}(P)$ itself, and from there the chain is the shared one: Theorem~\ref{thm:aba-hybrid-spec} and the safety transfer.
See \leanmodule{Leslie2Protocols/ABA/AFW/Composition} and \leanmodule{Leslie2Protocols/ABA/AFW/Substitution}.

\begin{definition}[The gather-based round families and their protocol-shaped systems]
  \label{def:aba-afw-round-families}
  \lean{PLTS.ABA.AFW.corruptionOverBracha}\lean{PLTS.ABA.AFW.corruptionOverBroadcastSpecification}\lean{PLTS.ABA.AFW.corruptionOverGatherSpecifications}\lean{PLTS.ABA.AFW.roundFamilyOverBracha}\lean{PLTS.ABA.AFW.roundFamilyOverBroadcastSpecification}\lean{PLTS.ABA.AFW.roundFamilyOverGatherSpecifications}\lean{PLTS.ABA.AFW.composed}\lean{PLTS.ABA.AFW.composedOverBroadcastSpecification}\lean{PLTS.ABA.AFW.composedOverGatherSpecifications}\leanok
  \uses{def:aba-round-over-bracha, def:aba-gbca-instance-by-abdy, def:aba-composed, def:idle-family, def:process-algebra, def:relabel}
  Each system over the two-gather round, over the extended alphabet it has itself, gathered into the $\mathbb{N}$-indexed family with the corruption broadcast, and put through the composed system's pipeline beside its other three components: the three systems $\afw{\mathrm{composed}}(P)$, $\afw{\mathrm{composed}}_{\mathrm{BRB}\,\mathrm{spec}}(P)$ and $\afw{\mathrm{composed}}_{\mathrm{Gather}\,\mathrm{spec}}(P)$, of which the first reads the gather-based protocol as a composition --- \leandecl{PLTS.ABA.AFW.composed}.
\end{definition}


\blueprintfig{fig-afw-composed}{The gather-based composed system \leandecl{PLTS.ABA.AFW.composed} and its three-stage substitution.
Component for component this is the composed system of Definition~\ref{def:aba-composed}: the same $n$ round loops, the same ABA network carrying the DECIDED sets and the corrupted set with its budget, the same lifted common coin, under the same two hidings.
What differs is the refinable position, and that it is replaced three times rather than once.
The position holds \leandecl{PLTS.ABA.AFW.roundFamilyOverBracha}, one composed round of Definition~\ref{def:aba-round-composition} per $r \in \mathbb{N}$, each carrying $4n + 2$ networks, and the substitution runs in three stages rather than one: the Bracha instances, then the gather instances, then the counting onto Transition System~2.
The intermediate systems are $\afw{\mathrm{composed}}_{\mathrm{BRB}\,\mathrm{spec}}(P)$ and $\afw{\mathrm{composed}}_{\mathrm{Gather}\,\mathrm{spec}}(P)$, and the third stage lands on $\mathrm{hybrid}(P)$ itself, where the two chains meet.
The other three components stand unchanged throughout.}
\begin{theorem}[The gather-based composition]
  \label{thm:aba-afw-protocol-composed}
  \lean{PLTS.ABA.AFW.roundProjection}\lean{PLTS.ABA.AFW.BoundInvariant}\lean{PLTS.ABA.AFW.ProtocolRelation}\lean{PLTS.ABA.AFW.protocolSimulation}\lean{PLTS.ABA.AFW.protocol_composed}\leanok
  \uses{def:aba-afw-protocol, def:aba-afw-round-families, thm:forward-sound, def:achievable}
  The gather-based protocol is forward-simulated by its composed system, along the Dirac lift of a relation that computes the composed state from the implementation's, \leandecl{PLTS.ABA.AFW.roundProjection}, and carries the bound bit of every round whose candidate is on record, \leandecl{PLTS.ABA.AFW.BoundInvariant}; hence $\operatorname{ATD}(\afw{\mathrm{protocol}}(P)) \subseteq \operatorname{ATD}(\afw{\mathrm{composed}}(P))$: \leandecl{PLTS.ABA.AFW.protocolSimulation}, \leandecl{PLTS.ABA.AFW.protocol_composed}.
  Each transition of the implementation is matched by one event of the composed round.
\end{theorem}

\begin{proof}\leanok
  \uses{def:aba-afw-protocol, def:aba-afw-round-families, def:forwardsim}
  See the Lean proof of \leandecl{PLTS.ABA.AFW.protocolSimulation}.
\end{proof}

\begin{theorem}[The gather-based substitution]
  \label{thm:aba-afw-substitution}
  \lean{PLTS.ABA.AFW.broadcastSubstitutionRelationFamily}\lean{PLTS.ABA.AFW.gatherSubstitutionRelationFamily}\lean{PLTS.ABA.AFW.roundSpecificationSubstitutionRelationFamily}\lean{PLTS.ABA.AFW.familyBroadcastSubstitutionSimulation}\lean{PLTS.ABA.AFW.familyGatherSubstitutionSimulation}\lean{PLTS.ABA.AFW.familyRoundSpecificationSubstitutionSimulation}\lean{PLTS.ABA.AFW.substitutionSimulation}\lean{PLTS.ABA.AFW.substitution}\leanok
  \uses{def:aba-afw-round-families, def:aba-hybrid, thm:gbca-round-substitutions, thm:gbca-round-refines-specification, thm:gbca-family-refines, thm:precong-parallel, thm:precong-abstract, thm:precong-relabel, thm:forward-sound}
  Each substitution inside the round, lifted by the family congruence and run under the composed system's own context through the four congruences, is a stage from one protocol-shaped system to the next; the third stage lands on $\mathrm{hybrid}(P)$ itself.
  Chaining the stages gives \leandecl{PLTS.ABA.AFW.substitutionSimulation} and the inclusion $\operatorname{ATD}(\afw{\mathrm{composed}}(P)) \subseteq \operatorname{ATD}(\mathrm{hybrid}(P))$: \leandecl{PLTS.ABA.AFW.substitution}.
\end{theorem}

\begin{proof}\leanok
  \uses{def:aba-afw-round-families, def:aba-hybrid, thm:gbca-round-substitutions, thm:gbca-round-refines-specification, thm:gbca-family-refines, thm:precong-parallel, thm:precong-abstract, thm:precong-relabel, thm:forward-sound}
  See the Lean proof of \leandecl{PLTS.ABA.AFW.substitution}.
\end{proof}


Below that point the chain has one step more than the direct one, the substitution splitting into its three stages:
\[ \begin{aligned}
  \afw{\mathrm{protocol}}(P)
  &\;\sqsubseteq\; \afw{\mathrm{composed}}(P)
  \;\sqsubseteq\; \afw{\mathrm{composed}}_{\mathrm{BRB}\,\mathrm{spec}}(P) \\
  &\;\sqsubseteq\; \afw{\mathrm{composed}}_{\mathrm{Gather}\,\mathrm{spec}}(P)
  \;\sqsubseteq\; \mathrm{hybrid}(P)
  \;\sqsubseteq\; \mathrm{spec}(P).
\end{aligned} \]
The first step is Theorem~\ref{thm:aba-afw-protocol-composed}, the middle three are the stages of Theorem~\ref{thm:aba-afw-substitution}, and the last is Theorem~\ref{thm:aba-hybrid-spec}, shared with the direct chain.
The chain is assembled in \leanmodule{Leslie2Protocols/ABA/Results}, beside the direct one.

\begin{theorem}[Safety of the gather-based protocol]
  \label{thm:aba-afw-main}
  \lean{PLTS.ABA.AFW.composed_refines}\lean{PLTS.ABA.AFW.composed_safe}\lean{PLTS.ABA.AFW.refines}\lean{PLTS.ABA.AFW.main}\lean{PLTS.ABA.AFW.chainSimulation}\leanok
  \uses{thm:aba-afw-protocol-composed, thm:aba-afw-substitution, thm:aba-hybrid-spec, thm:aba-spec-safe, thm:trace-support, thm:forward-trans}
  Every trace distribution achievable by the gather-based protocol as it runs is achievable by the ABA specification, and every positive-probability trace of it satisfies Validity and Agreement: \leandecl{PLTS.ABA.AFW.refines}, \leandecl{PLTS.ABA.AFW.main}; the composed simulation is \leandecl{PLTS.ABA.AFW.chainSimulation}.
  The same holds of the composed system alone: \leandecl{PLTS.ABA.AFW.composed_refines}, \leandecl{PLTS.ABA.AFW.composed_safe}.
\end{theorem}

\begin{proof}\leanok
  \uses{thm:aba-afw-protocol-composed, thm:aba-afw-substitution, thm:aba-hybrid-spec, thm:aba-spec-safe, thm:trace-support, thm:forward-trans}
  See the Lean proofs of \leandecl{PLTS.ABA.AFW.refines} and \leandecl{PLTS.ABA.AFW.main}.
\end{proof}


\section{The ghost-free system}

Both chains run over a system whose adversary keeps a ghost for each round, and that ghost is deleted here.
The route is the removal of Definition~\ref{def:auxiliary-variable-removal} at the network, carried through the composition pipeline by its congruences and collapsed at the hidden alphabet.
The construction is in \leanmodule{Leslie2Protocols/ABA/GhostRemoval/GhostFreeSystem}.
The two instantiations are in \leanmodule{Leslie2Protocols/ABA/GhostRemoval/ImplementationByABDY} and in \leanmodule{Leslie2Protocols/ABA/GhostRemoval/ImplementationByAFW}.

\begin{definition}[The ghost-free system]
  \label{def:aba-ghost-free}
  \lean{PLTS.ABA.Implementation.systemGhostFree}\lean{PLTS.ABA.ABDY.ghostFreeProtocol}\lean{PLTS.ABA.AFW.ghostFreeProtocol}\leanok
  \uses{def:aba-implementation, def:aba-protocol, def:aba-afw-protocol}
  The implementation of Definition~\ref{def:aba-implementation} at a one-element ghost record: the same programs, the same network transitions and the same coin oracle, with the adversary holding nothing of a round and its two graded-agreement returns free to announce either bit: \leandecl{PLTS.ABA.Implementation.systemGhostFree}.
  The ghost-free systems of the two protocols are that definition at the two implementations, \leandecl{PLTS.ABA.ABDY.ghostFreeProtocol} and \leandecl{PLTS.ABA.AFW.ghostFreeProtocol}, each the protocol of Definition~\ref{def:aba-protocol} and of Definition~\ref{def:aba-afw-protocol} with the record dropped and nothing else changed.
  The ghost output of D30 is a relation on the announced bit, which is what puts the three readings in one table: each algorithm instantiates it as the equation between the bit and the record it keeps, and the ghost-free system as the full relation.
\end{definition}

\begin{theorem}[The announced bit affects no trace distribution]
  \label{thm:aba-ghost-removal}
  \lean{PLTS.ABA.Implementation.system_ghostRemoval}\lean{PLTS.ABA.ABDY.protocol_ghostRemoval}\lean{PLTS.ABA.AFW.protocol_ghostRemoval}\lean{PLTS.ABA.ABDY.ghostFreeProtocol_safe}\lean{PLTS.ABA.AFW.ghostFreeProtocol_safe}\leanok
  \uses{def:auxiliary-variable-removal, def:aba-ghost-free, def:aba-labels, def:achievable, thm:aba-main, thm:aba-afw-main}
  Each protocol achieves exactly the trace distributions its ghost-free system achieves: \leandecl{PLTS.ABA.ABDY.protocol_ghostRemoval} and \leandecl{PLTS.ABA.AFW.protocol_ghostRemoval}, both instances of \leandecl{PLTS.ABA.Implementation.system_ghostRemoval}.
  No label map stands in either equation, a graded-agreement return being a label of the hidden alphabet of Definition~\ref{def:aba-labels}.
  Every conclusion of Theorem~\ref{thm:aba-main} and of Theorem~\ref{thm:aba-afw-main} therefore holds of the ghost-free system, Validity and Agreement included: \leandecl{PLTS.ABA.ABDY.ghostFreeProtocol_safe} and \leandecl{PLTS.ABA.AFW.ghostFreeProtocol_safe}.
\end{theorem}

\begin{proof}\leanok
  \uses{def:auxiliary-variable-removal, def:aba-ghost-free, def:aba-labels}
  See the Lean proof of \leandecl{PLTS.ABA.Implementation.system_ghostRemoval}.
  The two protocol statements are instances of it.
\end{proof}


\begin{thebibliography}{LeslieBP}

\bibitem[ABDY22]{ABDY22}
Ittai Abraham, Naama Ben-David and Sravya Yandamuri.
Efficient and adaptively secure asynchronous binary agreement via binding crusader agreement.
In \emph{41st ACM Symposium on Principles of Distributed Computing}, pages 381--391, 2022.

\bibitem[AFW25]{AFW25}
Hagit Attiya, Itay Flam and Jennifer L.\ Welch.
Why canonical rounds fail for optimal Byzantine resilience.
arXiv:2510.04310v2 [cs.DC], 2026.
Extended version of the paper in \emph{45th ACM Symposium on Principles of Distributed Computing}, pages 197--207, 2026;
citations here are to the extended version, whose appendices carry the gather algorithm and the approximate-agreement subroutine.

\bibitem[AL91]{AL91}
Mart\'in Abadi and Leslie Lamport.
The existence of refinement mappings.
\emph{Theoretical Computer Science}, 82(2):253--284, 1991.

\bibitem[Bra87]{Bra87}
Gabriel Bracha.
Asynchronous Byzantine agreement protocols.
\emph{Information and Computation}, 75(2):130--143, 1987.

\bibitem[LeslieBP]{LeslieBP}
Gaspard Reghem.
\emph{Verifying ABA with Leslie}.
Project blueprint, 2026. The source document whose Transition Systems and Algorithms the encodings of this section follow.

\bibitem[Segala95]{Segala95}
Roberto Segala.
A compositional trace-based semantics for probabilistic automata.
In \emph{CONCUR~'95}, pages 234--248. Springer, 1995.

\end{thebibliography}
